\documentclass[a4paper,11pt]{amsart}
\pdfoutput=1   %%pro arXiv
\usepackage[utf8]{inputenc}
\usepackage[english]{babel}
\usepackage{indentfirst}
\usepackage{csquotes}
\usepackage{soul} 
\sethlcolor{green} 
\usepackage{enumerate}

\sloppy 

\usepackage{amsmath,amsthm,amssymb,mathrsfs,thmtools}
\usepackage{dsfont}

\usepackage{enumitem}
\usepackage[toc,page]{appendix}

\usepackage{graphicx}
\usepackage{mathtools,array}
\usepackage[all]{xy}
\usepackage{tikz}
\usetikzlibrary{cd}
\usetikzlibrary{fit}
\usetikzlibrary{knots}
\usetikzlibrary{positioning}
\usetikzlibrary{arrows.meta}

\setlength{\headheight}{12pt}% ...at least 51.60004pt

\usepackage{fancyhdr,floatpag}
\usepackage{upgreek,color}
\usepackage{hyperref}
%\hypersetup{colorlinks=false,bookmarks=true}

\usepackage{blkarray}
\usepackage{multirow}

\usepackage{dutchcal} %calligraphic letters in math mode

\usepackage[style=numeric,sorting=nyt,maxbibnames=7,maxcitenames=5,block=space,backend=bibtex]{biblatex}
\addbibresource{biblio.bib}
% \usepackage{makeidx}
% \makeindex
\renewbibmacro{in:}{}

\pagestyle{fancy}
\fancyhead{} 
\fancyfoot{} 
\fancyfoot[LE,RO]{\thepage}
\fancyhead[RO]{Reidemeister numbers for soluble arithmetic groups}
\fancyhead[LE]{P. Lins and Y. Santos}
\renewcommand{\headrulewidth}{0.4pt}
\renewcommand{\footrulewidth}{0pt}
\renewcommand{\footruleskip}{-0.5cm}

\hypersetup{linkbordercolor={0 0.60 0.75}, linkcolor=MidnightBlue}
\hypersetup{citebordercolor={0 0.60 0.75}, citecolor=MidnightBlue}

\renewcommand{\theenumi}{{{\roman{enumi}}}}
\renewcommand{\theenumii}{\alph{enumii}}

\numberwithin{equation}{section}
\theoremstyle{plain}
\newtheorem{theorem}{Theorem}
\newtheorem{thm}{Theorem}[section]
\newtheorem{con}[thm]{Conjecture}
\newtheorem{conjecture}[theorem]{Conjecture}
\newtheorem{cor}[thm]{Corollary}
\newtheorem{corollary}[theorem]{Corollary}
\newtheorem{dfn}[thm]{Definition}
\newtheorem{lem}[thm]{Lemma}
\newtheorem{obs}[thm]{Observation}
\newtheorem{pps}[thm]{Proposition}
\newtheorem{cla}[thm]{Claim}

\theoremstyle{definition}
\newtheorem{rmkx}[thm]{Remark}
\newenvironment{rmk}
  {\pushQED{\qed}\renewcommand{\qedsymbol}{$\triangle$}\rmkx}
  {\popQED\endrmkx}
\newtheorem{examplex}[thm]{Example}
\newenvironment{exm}
  {\pushQED{\qed}\renewcommand{\qedsymbol}{$\triangle$}\examplex}
  {\popQED\endexamplex}
\newtheorem{prob}[thm]{Problem}
\newtheorem{qst}[thm]{Question}

\declaretheoremstyle[
  spaceabove=-6pt,
  spacebelow=6pt,
  headfont=\normalfont\bfseries,
  postheadspace=1em,
  qed=\qedsymbol,
  headpunct={}
]{mystyle} 
\declaretheorem[name={proofof},style=mystyle,unnumbered,
]{proofof}  
\makeatletter
\renewenvironment{proofof}[1] {\par\pushQED{\qed}\normalfont\topsep6\p@\@plus6\p@\relax\trivlist  \item[\hskip\labelsep
        \bfseries
    Proof of #1.]\ignorespaces}{\popQED\endtrivlist\@endpefalse}
\makeatother
\makeatletter
\newcommand{\thickhline}{
    \noalign {\ifnum 0=`}\fi \hrule height 1pt
    \futurelet \reserved@a \@xhline
}
\newcolumntype{"}{@{\hskip\tabcolsep\vrule width 1pt\hskip\tabcolsep}}
\makeatother

\makeatletter
\renewenvironment{proof}[1][\proofname] {\par\pushQED{\qed}\normalfont\topsep6\p@\@plus6\p@\relax\trivlist\item[\hskip\labelsep\bfseries#1\@addpunct{.}]\ignorespaces}{\popQED\endtrivlist\@endpefalse}
\makeatother

\newcommand\Tstrut{\rule{0pt}{2.6ex}}         
\newcommand\Bstrut{\rule[-0.9ex]{0pt}{0pt}}  

\makeatletter
\renewcommand{\@secnumfont}{\bfseries}
\makeatother
\usepackage{etoolbox}
\makeatletter
\patchcmd{\section}{\scshape}{\bf}{}{}
\patchcmd{\subsection}{\scshape}{\bf}{}{}
\patchcmd{\subsubsection}{\scshape}{\bf}{}{}
\makeatother

%%% Numbers and rings
\newcommand{\N}{\mathbb{N}} 
\newcommand{\Z}{\mathbb{Z}} 
\newcommand{\Zum}{\mathbb{Z}_{\geq 1}} 
\newcommand{\Nzero}{\mathbb{Z}_{\geq 0}} 
\newcommand{\Q}{\mathbb{Q}} 
\newcommand{\R}{\mathbb{R}} 
\newcommand{\K}{\mathbb{K}} 
\newcommand{\F}{\mathbb{F}} 
\newcommand{\C}{\mathbb{C}} 
\newcommand{\real}{\mathfrak{Re}} 
\newcommand{\imag}{\mathfrak{Im}} 
\newcommand{\pai}{\mathbb{Z}_{p}} 
\newcommand{\laurent}[1]{#1[t,t^{-1}]} 
\newcommand{\OS}{\mathcal{O}_S} 

\newcommand{\anel}[1]{#1_{\text{ring}}} %subscript 'ring', e.g. for Aut_ring(R)
\newcommand{\addi}[1]{#1_{\text{add}}} %subscript 'add', which stands for 'additive', e.g. for \sigma_add

\newcommand{\p}{\mathfrak{p}} 
\newcommand{\Pp}{\mathfrak{P}}
\newcommand{\ri}{\mathcal{O}} 
\newcommand{\ris}{\mathcal{O}_S} %ring of S-integers 2
\newcommand{\lri}{\mathfrak{o}}
\newcommand{\Lri}{\mathfrak{O}}

\newcommand{\riK}{\mathcal{O}_{\mathbb{K}}} %ring of integers of a number field K

%%% Group schemes and matrices
\DeclareMathOperator{\GL}{GL} %general linear
\DeclareMathOperator{\GE}{GE} %elementary general linear
\DeclareMathOperator{\SL}{SL} %special linear
\DeclareMathOperator{\Sp}{Sp} %symplectic
\DeclareMathOperator{\SO}{SO} %special orthogonal
\newcommand{\Eij}{{\rm E}_{ij}} %elementary matrix in the Lie algebra ('zeroes elsewhere')
\newcommand{\eij}{e_{i,j}} %elementary matrix in the group
\newcommand{\di}{d_i} %elementary diagonal matrix in the group
\newcommand{\dk}[1]{d_{#1}} %elementary diagonal matrix in the group with index-inputs
\newcommand{\Ekl}[1]{{\rm E}_{#1}} %elementary matrix in the Lie algebra with index-inputs
\newcommand{\ekl}[1]{e_{#1}} %elementary matrix in the group with index-inputs
\newcommand{\sbgpeij}{\mathcal{E}_{i,j}} %subgroup generated by the eij(r)
\newcommand{\sbgpekl}[1]{\mathcal{E}_{#1}} %subgroup generated by the eij(r), with index-inputs
\newcommand{\mbU}{\mathbf{U}} %unipotent (or upper unitriangular in GLn)
\newcommand{\mbD}{\mathbf{D}} %diagonal (or torus in GLn)
\DeclareMathOperator{\Diag}{Diag} %diagonal matrices
\newcommand{\sbgpdi}{\mathcal{D}_{i}} %subgroup generated by the d_i(u)
\newcommand{\sbgpdk}[1]{\mathcal{D}_{#1}} %subgroup generated by the d_i(u), with index-input
\newcommand{\projective}[1]{\mathbb{P}{#1}}
\newcommand{\PGL}{\mathbb{P}\mathrm{GL}}
\newcommand{\PSL}{\mathbb{P}\mathrm{SL}}
\newcommand{\Mult}{\mathbb{G}_{m}} %multiplicative group scheme
\newcommand{\Addi}{\mathbb{G}_{a}} %additive group scheme
\newcommand{\Aff}{\mathbb{A}\mathrm{ff}} %the affine group
\newcommand{\AffR}[1]{\mathbb{G}_{a}(#1) \rtimes \mathbb{G}_{m}(#1)} %affine group with input
\newcommand{\Bzero}{\mathbf{B}_{2}^{\circ}} %upper triangular in SL2
\newcommand{\BzeroR}{\mathbf{B}_{2}^{\circ}(R)} %same, over a ring R
\newcommand{\mbP}{\mathbf{P}} %parabolic
\newcommand{\mbB}{\mathbf{B}} %Borel (or upper triangular in GLn)
\newcommand{\mbG}{\mathbf{G}} %'generic' group scheme
\newcommand{\mbH}{\mathbf{H}} %'generic' or torus
\newcommand{\mbT}{\mathbf{T}} %torus
\newcommand{\mbL}{\mathbf{L}} %Levi
\newcommand{\mbW}{\mathbf{W}} %Z sd D
\newcommand{\PB}{\mathbb{P}\mathbf{B}} %projective Borel

\newcommand{\PBplus}{\mathbb{P}B^{+}} %projective version of B_n^+, without the index n

%%% Better symbols and tools
\newcommand{\phee}{\varphi} %best phi
\newcommand{\veps}{\varepsilon} %best epsilon
\newcommand{\cupdot}{\mathbin{\mathaccent\cdot\cup}} %disjoint union
\newcommand{\barra}[1]{\overline{#1}}
\newcommand{\chapeu}[1]{\widehat{#1}}
\newcommand{\til}[1]{\widetilde{#1}}
\newcommand{\gera}[1]{\langle {#1} \rangle}
\newcommand{\set}[1]{\{ #1 \}}
\newcommand{\vazio}{\varnothing} %empty set
\newcommand{\lbr}{\llbracket}
\newcommand{\rbr}{\rrbracket}
\newcommand{\mbb}{\mathbb}
\newcommand{\mb}{\mathbf}
\newcommand{\mc}{\mathcal}
\newcommand{\mf}{\mathfrak}
\newcommand{\mscr}{\mathscr}
\newcommand{\mt}[1]{\text{#1}}
\newcommand{\into}{\hookrightarrow}
\newcommand{\onto}{\twoheadrightarrow}
\newcommand{\nsgp}{\trianglelefteq} 
\newcommand{\cs}{\leq_c} 
\newcommand{\os}{\leq_o} 
\newcommand{\cns}{\triangleleft_c} 
\newcommand{\ons}{\triangleleft_o} 

\newcommand{\Ri}{R_{\infty}} % R_infty
\newcommand{\propRi}{property~$\Ri$} %property R_infty
\newcommand{\PropRi}{Property~$\Ri$} %Property R_infty
\newcommand{\RSpec}{\mathrm{RSpec}}
\newcommand{\bpt}{B_{2}^{+}(\ris)} %"'B-plus-two"'
\newcommand{\bt}{B_{2}(\ris)} %"'B-two"'
\newcommand{\bpn}{\mathbf{B}_{n}^{+}(\ris)}
\newcommand{\bnr}{\mathbf{B}_{n}(R)}
\newcommand{\dpt}{\mathbf{D}_{2}^{+}(\ris)}
\newcommand{\dpn}{\mathbf{D}_{n}^{+}(\ris)}
\newcommand{\ut}{U_{2}(\ris)}
\newcommand{\un}{U_{n}(\ris)}
\newcommand{\e}{\mathbf{e}}

\newcommand{\Fn}[1]{\mathrm{F}_{#1}} %finiteness properties
\newcommand{\FPn}[1]{\mathrm{FP}_{#1}}
\newcommand{\bref}[1]{{\bf(\ref{#1})}} %bracket reference
\newcommand{\dref}[2]{\ref{#1}{\bf(\ref{#2})}} %double reference, e.g. Lemma 1(ii)



%%% Operators
\DeclareMathOperator{\Mat}{Mat}
\DeclareMathOperator{\Hom}{Hom}
\DeclareMathOperator{\Det}{Det}
\DeclareMathOperator{\End}{End}
\DeclareMathOperator{\Aut}{Aut}
\DeclareMathOperator{\Inn}{Inn}
\DeclareMathOperator{\im}{Im}
\DeclareMathOperator{\rk}{rk}
\DeclareMathOperator{\vspan}{span}
\DeclareMathOperator{\Obj}{Obj}
\DeclareMathOperator{\CAT}{CAT}
\DeclareMathOperator{\Lie}{Lie}
\DeclareMathOperator{\ad}{ad}
\DeclareMathOperator{\carac}{char}
\DeclareMathOperator{\Tor}{Tor} 
\DeclareMathOperator{\Frac}{Frac} 
\DeclareMathOperator{\id}{id}

\renewcommand*{\bibfont}{\small} 

\title[Reidemeister numbers for soluble arithmetic groups]{Reidemeister numbers for arithmetic Borel subgroups in type A}
\author{Paula Macedo Lins de Araujo and Yuri Santos Rego}
\address{Katholieke Universiteit Leuven, \newline 
Wiskunde, Campus Kulak, \newline 
Etienne Sabbelaan 53, bus 7657, \newline 
8500 Kortrijk, Belgi\"e}
\curraddr{University of Lincoln, \newline 
Charlotte Scott Centre for Algebra, School of Mathematics and Physics, \newline 
Isaac Newton Building, Brayford Pool, LN6 7TS, Lincoln, United Kingdom}
\email{pmacedolinsdearaujo@lincoln.ac.uk}
\address{Otto-von-Guericke-Universit\"at Magdeburg, \newline 
Fakult\"at f\"ur Mathematik -- Institut f\"ur Algebra und Geometrie, \newline 
Postfach 4120, 39016 Magdeburg, Deutschland}
\email{yuri.santos@ovgu.de}
\subjclass[2020]{20E36, 20F16, 20G30}
%; 11C99, 13G05}
\keywords{Reidemeister numbers, soluble groups, $S$-arithmetic groups, property $R_\infty$, upper triangular matrices.}

\begin{document}
\begin{abstract}
The Reidemeister number $R(\varphi)$ of a group automorphism $\varphi \in \mathrm{Aut}(G)$ encodes the number of orbits of the $\varphi$-twisted conjugation action of $G$ on itself, and the Reidemeister spectrum of $G$ is defined as the set of Reidemeister numbers of all of its automorphisms. We obtain a sufficient criterion for some groups of triangular matrices over integral domains to have property~$R_\infty$, which means that their Reidemeister spectrum equals $\{\infty\}$. Using this criterion, we show that Reidemeister numbers for certain soluble $S$-arithmetic groups behave differently from their linear algebraic counterparts --- contrasting with results of Steinberg, Bhunia, and Bose.  
\end{abstract}
\maketitle \vspace{-1.0cm}
\thispagestyle{empty}

\section{Introduction} 
Given a group~$G$ and an automorphism $\phee\in\Aut(G)$, the $\phee$\emph{-(twisted) conjugacy class} of $g\in G$ is the set of elements that are $\phee$-conjugate to $g$, 
\[ [g]_\phee = \set{h g \phee(h)^{-1} \mid h \in G}. \]
The \emph{Reidemeister number} of $\phee$, denoted by $R(\phee)$, is the total number of $\phee$-conjugacy classes (also called {Reidemeister classes} of $\phee$). In particular, the \emph{class number} of a group is just $R(\mathrm{id})$. 
The \emph{Reidemeister spectrum} $\RSpec(G)$ of the group $G$ is defined as the set 
\[\RSpec(G)= \{n \in \Z_{\geq 1} \cup \{\infty\} \mid \exists \phee \in \Aut(G) \text{ with } R(\phee)=n\}\] 
of all possible Reidemeister numbers for $G$. Following Taback and Wong~\cite{TabackWong0}, we say that $G$ has \emph{property}~$\Ri$ when $\RSpec(G) = \{\infty\}$.

The main interest in a group having or not~$\Ri$ is due to connections with fixed point statements in different areas; we refer the reader, e.g., to \cite[Theorem~6.1]{DacibergWongCrelle}, \cite[Theorem~10.1]{SteinbergEndo}, and \cite[Theorem~5.4]{Brasuca0} for examples in algebraic topology, Lie theory, and cohomology, respectively. 

For a connected linear algebraic group $\mbB$ defined over an algebraically closed field, there is a strong relation between its structure and its Reidemeister spectrum, as investigated by Steinberg~\cite{SteinbergEndo} and more recently by Bhunia and Bose~\cite{BhuniaBose1,BhuniaBose2}. Indeed, if there exists some \emph{algebraic} automorphism $\phee$ of $\mbB$ with $R(\phee) < \infty$, then $\mbB$ is necessarily soluble~\cite[Theorem~17]{BhuniaBose1}. As a partial converse, if $\mbB$ is soluble then $R(\phee)$ is either $1$ or $\infty$~\cite[Theorem~2.4]{BhuniaBose2}, and the latter always holds in case $\mbB \leq \mbG$ is a Borel subgroup of a semisimple group $\mbG$~\cite[Theorem~2.10]{BhuniaBose2}. 

Turning to abstract countable groups, recall that linear algebraic groups have prominent `discrete counterparts'. Let $\mbB$ be a linear algebraic group defined over a global field $\K$. A group $\Gamma$ is called an $S$-arithmetic subgroup of $\mbB$ if $\Gamma$ is commensurable with the subgroup $\rho^{-1}(\GL_n(\OS)) \leq \mbB(\K)$ for a ring of $S$-integers $\OS \subset \K$ and some faithful $\K$-representation $\rho : \mbB \into \GL_n$. 
The reader is referred to the classics \cite{Margulis,PlatonovRapinchuk} for major results and literature on such groups. We briefly mention that the family of $S$-arithmetic groups includes important familiar objects, such as finitely generated free groups, crystallographic groups, and many lattices in products of (real, complex, and $p$-adic) semisimple Lie groups.

The results of Steinberg and Bhunia--Bose mentioned above have natural follow-up questions in the arithmetic set-up. Firstly, how does the Reidemeister spectrum of an $S$-arithmetic group affect its (algebraic) structure? Secondly, if $\Gamma$ is an $S$-arithmetic subgroup of $\mbB$, with $\mbB$ a Borel subgroup of a semisimple group over the global field $\K$, does $\Gamma$ necessarily have property~$\Ri$? Lastly, in case such a $\Gamma$ \emph{does not} have~$\Ri$, does it hold $\RSpec(\Gamma) = \{1, \infty\}$? 

A partial answer to the first question follows from work of Jabara. Namely, if a finitely generated linear group $\Gamma$ has an automorphism $\phee$ of \emph{finite} order with $R(\phee) < \infty$, then $\Gamma$ is indeed virtually soluble; see \cite[Theorem~C]{Jabara}. Our contribution to the theory is the following result, answering the other two questions in the negative --- although barely, since~$\Ri$ actually \emph{is} achieved in higher solubility classes, in analogy to known cases of nilpotent groups~\cite{DacibergWongCrelle,Romankov,KarelDacibergNil,TimurUniTri}.

\begin{thm} \label{thm:Lieapplication}
Let $p \in \Z_{\geq 0}$ be zero or an arbitrary prime number, and write $\PB_n$ for the Borel subgroup of (upper) triangular matrices of the split semisimple group $\PGL_n$. Then for every $n \geq 2$ there is a global field $\K_{n,p}$ of chosen characteristic $p \geq 0$ and a family of (\emph{non-nilpotent}, class-$n$ soluble) $S$-arithmetic subgroups $\Gamma_{n,p}$ of $\PB_n$ over $\K_{n,p}$ such that 
\begin{enumerate}
    \item $\RSpec(\Gamma_{2,p}) \neq \{1,\infty\}$, but 
    \item $\RSpec(\Gamma_{n,p}) = \{\infty\}$ for all $n \geq 4$.
\end{enumerate}
In case $\mathrm{char}(\K_{n,p}) = p > 0$, one can take the groups $\Gamma_{n,p}$ to be all finitely generated, or all non-finitely generated. (If $\mathrm{char}(\K_{n,p}) = p = 0$, all such $\Gamma_{n,p}$ are finitely presented.)
\end{thm}

It should be stressed that Theorem~\ref{thm:Lieapplication} and the previously mentioned results are `uniform' in the sense that they hold in every characteristic. This is in strong contrast with other results for algebraic or $S$-arithmetic groups; cf.~\cite{LitterickThomasCRgood,Tiemeyer,Bux04}, for example. 

In characteristic zero, a weaker form of Theorem~\ref{thm:Lieapplication} can be deduced from Nasybullov's work~\cite{TimurUniTri} by taking a family of nilpotent arithmetic groups. Requiring non-nilpotent examples and achieving~$\Ri$ in class $4$ onwards, as in our theorem, improves on Nasybullov's findings. In positive characteristic, examples of metabelian $S$-arithmetic groups as in part~(i) of Theorem~\ref{thm:Lieapplication} were known by work of Gon\c{c}alves--Wong~\cite{DacibergWongWreath}. The case of arbitrarily high solubility class in positive characteristic is, to our knowledge, new. 

The main technical contribution of the present paper, which yields part~(ii) of Theorem~\ref{thm:Lieapplication},  
is a tool to check for property~$\Ri$ in linear groups of triangular matrices.

To be precise, let $\mbB_n \leq \GL_n$ be the $\Z$-subscheme of (upper) triangular matrices. Given a (commutative, unital) ring $R$, we work concretely with the groups of $R$-points 
\[ 
\mbB_n(R) = \left( \begin{smallmatrix} * & * & \cdots & * \\  & * & & \\ & & \ddots & \vdots \\ & & & * \\  \end{smallmatrix} \right) \leq \GL_n(R).
\]
Additionally, we consider the projective variant $\PB_n(R)$ of $\mbB_n(R)$ obtained by modding out scalar matrices. Our main result is the following.

\begin{thm} \label{thm:Additive}
Suppose $R$ is an integral domain with finitely generated group of units. Given a ring automorphism $\alpha \in \anel{\Aut}(R)$, write $\addi{\alpha}$ for the same map seen as an automorphism of $(R,+)$, and let $\tau_\alpha \in \Aut( R\times R)$ be the induced `flip' automorphism $\tau_\alpha((r,s)) = (\addi{\alpha}(s),\addi{\alpha(r)})$. If $R(\addi{\alpha}) = \infty = R(\tau_\alpha)$ for every $\alpha \in \anel{\Aut}(R)$, then $G(R)$ has~$\Ri$ for all $G \in \{ \mbB_n, \PB_n \mid n \geq 4 \}$.
\end{thm}

We also give examples of \emph{non}-$S$-arithmetic groups to which Theorem~\ref{thm:Additive} applies; cf. Section~\ref{applications}. Putting our results further into perspective, Theorem~\ref{thm:Additive} is a contribution towards the wide problem of classifying which amenable linear groups have property~$\Ri$. There has been substantial progress in the (virtually) nilpotent \cite{DacibergWongCrelle,Romankov,KarelDacibergNil,TimurUniTri}, polycyclic \cite{KarelPenni,KarelSam,KarelSamIrisSolv}, and metabelian cases \cite{TabackWong0,DacibergWongWreath,FelshtynDacibergMetabelian}. Moreover, nonpolycyclic nilpotent-by-abelian groups of type~$FP_\infty$ always have~$\Ri$; see \cite[Theorem~4.3]{DesiDaciberg}. 

While Theorem~\ref{thm:Additive} has intersections with some of the above mentioned results, it contributes towards a systematic investigation of many soluble linear groups and elucidates how the structure of the underlying base ring leads to the given group having~$\Ri$. In a companion paper~\cite{Bn1}, whose first draft included most of the present article, we further highlight how the structure of the ground ring---in that case, as a module over its units---can detect~$\Ri$ for $\mbB_n(R)$ and its variants, starting from solubility class $n = 2$. In particular, it is possible to construct (in any characteristic $p \geq 0$) a family $\{\Gamma_{n,p}\}_{n \geq 2}$ of finitely \emph{presented}, class-$n$ soluble, non-nilpotent $S$-arithmetic groups $\Gamma_{n,p}$ with $\RSpec(\Gamma_{n,p}) = \{\infty\}$ for all $n$ and $p$.

Although we explicitly compute Reidemeister numbers in some cases (see Section~\ref{sec:newexampleswithoutRinfty}), we do not go so far as to determine the full Reiedemeister spectrum of the groups without~$\Ri$ that we investigate; see \cite{KarelPenni,FelshtynDacibergMetabelian,KarelKaiserSam,KarelSamIrisSolv} for examples in this direction. Similarly, we do not have descriptions for Reidemeister classes or asymptotic properties of Reidemeister numbers. These are interesting (and challenging) problems on their own. Related topics include problems around Higman's conjecture~\cite{HalasiPalfy,Higman,VLA}, dynamical zeta functions attached to Reidemeister numbers~\cite{DekTerVandeBus,FelshtynHill,FelshtynZietek}, and zeta functions of groups counting Reidemeister classes~\cite{PaulaII,ask,RoVo19,duSau05}. Lastly, we do not cover the other classical Lie types $\mathtt{B}_n$ to $\mathtt{G}_2$. Our results and the indicated literature thus motivate us to pose the following.

\begin{prob}
Let $\mc{G}$ be a split reductive linear algebraic group and fix a Borel subgroup $\mc{B} \leq \mc{G}$, all defined over some global field $\K$. Given an $S$-arithmetic subgroup $\Gamma$ of $\mc{B}$, describe the Reidemeister spectrum $\RSpec(\Gamma)$. And, in case $\RSpec(\Gamma) \neq \{\infty\}$, how do the Reidemeister numbers of automorphisms of $\Gamma$ grow?
\end{prob}

This paper is structured as follows. 
We start with Section~\ref{sec:LemmataRinfty} collecting well-known facts about Reidemeister numbers to be used throughout. In the same section we introduce our groups of interest in more detail. 
An important ingredient in the proof of Theorem~\ref{thm:Additive} is a description of automorphisms of the group of unitriangular matrices, due to V.~Lev\v{c}uk~\cite{LevchukOriginal} --- we present his theorem in our terminology in Section~\ref{sec:Levchuk}, and sketch how to deduce it from~\cite{LevchukOriginal}. 

We restate and prove Theorem~\ref{thm:Additive} in Section~\ref{sec:provadothmB}. In Section~\ref{applications}, we show how to apply it via nontrivial examples; see Proposition~\ref{pps:FptFptt-1OK}. 
Following up, we complement the previous findings with examples of groups \emph{without} property~$\Ri$, computing explicit Reidemeister numbers; cf. Proposition~\ref{nori}.  
For completion, the final Section~\ref{sec:backtoarithmetic} contains a proof of Theorem~\ref{thm:Lieapplication}. 

To keep the article as self-contained as possible 
we do not assume familiarity with 
algebraic or $S$-arithmetic groups and $S$-arithmetic rings.  
We work concretely with the given matrix groups and base rings, so the reader familiar with standard results might want to skip directly to the the proof of the main theorem in Section~\ref{sec:provadothmB}, and the applications in Section~\ref{applications}.

%%%
\section{Auxiliary results, and structure of the groups}
\label{sec:LemmataRinfty}
We first recall known results on Reidemeister numbers, particularly for group extensions. Then, we collect some facts and notation concerning the groups considered in this paper.

\begin{lem}[{See \cite[Cor.~2.5]{FelshtynTroitskyCrelle}}] \label{lem:ignoreinner}
Let $\phee \in \Aut(G)$ and let $\iota_g \in \mathrm{Inn}(G) \leq \Aut(G)$ denote the inner automorphism $\iota_g(h) = ghg^{-1}$. Then $R(\iota_g \circ \phee) = R(\phee)$.
\end{lem}

\begin{lem}[{See \cite[Prop.~1.2]{Daciberg}, \cite[Thm.~1]{WongCrelle}, \cite[Lem.~1.1(2)]{DacibergWongCrelle}}] \label{lem:desempre}
Suppose there is a short exact sequence $N \into G \onto Q$ where $N$ is invariant under $\phee\in\Aut(G)$. Denote by $\phee' \in \Aut(N)$ and $\overline{\phee} \in \Aut(Q)$ the automorphisms induced by $\phee$, i.e., $\phee' = \phee|_N$ and $\barra{\phee}(gN) := \phee(g)N$. Then the following hold.
\begin{enumerate}
\item \label{PHEEvsPHEEbarra} $R(\phee) \geq R(\overline{\phee})$.
\item \label{PHEEbarra_and_innerPHEE'} 
If $R(\overline{\phee}) < \infty$ and if $R(\iota_g \circ \phee') < \infty$ for all inner automorphisms $\iota_g \in \Aut(G)$, 
then $R(\phee) < \infty$. 
\item \label{PHEEbarrafinite_and_PHEE'infinite} Suppose $\overline{\phee}$ has \emph{finitely many} fixed points. If $R(\overline{\phee}) < \infty$ and $R(\phee') = \infty$, then $R(\phee) = \infty$.
\end{enumerate}
\end{lem}

\begin{lem}[{See \cite[Section~2 and Lemma~4.1]{Romankov}, \cite{KarelDacibergAbelian}}] \label{lem:ReidemeisterAbelian}
Let $A$ be a \emph{finitely generated} abelian group and $\phee \in \Aut(A)$. 
Then $R(\phee) = \infty$ if and only if $\phee$ has \emph{infinitely many} fixed points.
\end{lem}

We now turn to 
our groups of interest. 
For us, $R$ denotes a \emph{commutative} ring \emph{with unity} unless stated otherwise. 

The group $\mbB_n(R)$, also called the (standard) Borel subgroup of $\GL_n(R)$, is the soluble subgroup of $\GL_n(R)$ of upper triangular matrices. 

Let $\mbU_n(R) \leq \mbB_n(R)$ be the subgroup of unitriangular matrices, and $\mbD_n(R) \leq \mbB_n(R)$ 
the subgroup of diagonal matrices. As an abstract group, $\mbB_n(R)$ decomposes as the semi-direct product $\mbB_n(R) = \mbU_n(R) \rtimes \mbD_n(R)$. 
Since $R$ is commutative, the group $\mbU_n(R)$ is nilpotent (of class $n-1$) and {$\mbD_n(R)$} %\cong \Mult(R)^n$ 
is abelian, hence $\mbB_n(R)$ is soluble of class at most $n$.

Modding out scalar matrices leads to an important quotient. Specifically, 
consider the central subgroup 
\[Z_n(R) = \set{u\cdot \mb{1}_n \in \GL_n(R) \mid u \in R^\times},\] 
where $\mb{1}_n$ denotes the $n \times n$ identity matrix.  
The \emph{projective upper triangular group} $\PB_n(R)$ 
is given by 
\[ \PB_n(R) = \frac{\mbB_n(R)}{Z_n(R)} = \mbU_n(R) \rtimes \frac{\mbD_n(R)}{Z_n(R)}.\] 

\begin{rmk} \label{obs:centerBn} If $R$ is an integral domain, then $Z_n(R)$ 
is either trivial or coincides with the center of $\mbB_n(R)$. 
Thus $\PB_n(R)$ 
is a quotient of $\mbB_n(R)$ 
modulo a \emph{characteristic} subgroup whenever $R$ is an integral domain. 
\end{rmk}

We remark that $\PB_2(R)$ is isomorphic to the group $\Aff(R) \cong \left( \begin{smallmatrix} * & * \\ 0 & 1 \end{smallmatrix} \right) \leq \GL_2(R)$ of affine transformations of the base ring $R$; cf. \cite[Section~3]{Bn1}.

\medskip

We introduce some more notation that will be used throughout. We let $\Addi(R) = (R,+)$ and $\Mult(R) = (R^\times, \cdot)$ denote the underlying additive group and the group of units of the ring $R$, respectively.

For $r \in R$ and $1 \leq i \neq j \leq n$, we write $\eij(r) \in \GL_n(R)$ for the elementary matrix whose off-diagonal $(i,j)$-entry is $r \in R$, all its diagonal entries equal~$1$, and having zeroes elsewhere. For instance, in $\GL_2(\Z)$ one has $\ekl{1,2}(2) = \left( \begin{smallmatrix} 1 & 2 \\ 0 & 1 \end{smallmatrix} \right)$. For a fixed $(i,j)$ we also denote $\mc{E}_{i,j} := \langle \{ \ekl{i,j}(r) \in \mbU_n(R) \mid r \in R \}\rangle \cong \Addi(R)$.

Given $i \in \set{1,\ldots,n}$ and $u \in R^\times$ we let $\di(u) \in \GL_n(R)$ denote the diagonal matrix whose $i$-th entry is $u$ and all other (diagonal) entries equal~$1$. For example, $\dk{2}(-1) = \left( \begin{smallmatrix} 1 & 0 \\ 0 & -1 \end{smallmatrix} \right)$ in $\GL_2(\Z)$. 
The elements 
$\eij(r)$ are called \emph{elementary} matrices and the $\di(u)$ are 
\emph{elementary diagonal} matrices.

We recall that elements of $\mbU_n(R)$ can be uniquely written as a product of elementary matrices ordered according to the superdiagonals of $\mbU_n(R)$. More precisely, given $x \in \mbU_n(R)$, there exist (uniquely determined) $r_{i,j} \in R$ such that 
\begin{align} \label{rel:elementsofUn} 
\begin{split}
x = & \phantom{.} e_{1,2}(r_{1,2}) e_{2,3}(r_{2,3}) \cdot \ldots \cdot e_{n-1,n}(r_{n-1,n}) \cdot \\
 & \cdot  e_{1,3}(r_{1,3}) \cdot \ldots \cdot e_{n-2,n}(r_{n-2,n}) \cdot \ldots \ldots \cdot e_{1,n}(r_{1,n}). 
\end{split}
\end{align}

All of the group theoretical properties above for $\mbB_n(R)$, $\mbD_n(R)$ and $\mbU_n(R)$ are deduced from classic observations using the following sets of relations; see, e.g., \cite{HahnO'Meara, Silvester}. 
(Commutators in this paper are written as $[g,h] = ghg^{-1}h^{-1}$.) 
\begin{align}
\label{rel:commutators}
\begin{split}
[\eij(r),\ekl{k,l}(s)] & =
\begin{cases}
\ekl{i,l}(rs) & \mbox{if } j=k,\\
1 & \mbox{if } i \neq l \text{ and } k \neq j,
\end{cases} \\
\di(u) \ekl{k,l}(r) \di(u)^{-1} & = 
\begin{cases}
\ekl{k,l}(ur) & \mbox{if } i=k,\\
\ekl{k,l}(u^{-1}r) & \mbox{if } i=l, \\
1 & \mbox{otherwise}.
\end{cases}
\end{split}
\end{align}
The equations~\bref{rel:commutators} are also referred to as \emph{elementary} or \emph{commutator relations}.

Elements of the diagonal part of $\PB_n(R)$ 
shall be denoted by $[d]$. 
Thus, any element $g \in \PB_n(R)$ 
can be uniquely written as $g = u[d]$ where $u \in \mbU_n(R)$ and $d \in \mbD_n(R)$. 
Moreover, the elementary relations~\bref{rel:commutators} yield 
\begin{align} \label{rel:projectiveconjugation}
[d] \ekl{k,l}(r) [d]^{-1} = \ekl{k,l}(u_k u_l^{-1} r) & \phantom{a} & \text{ if } d = d_1(u_1) \cdots d_n(u_n).
\end{align}

In the beginning of Section~\ref{sec:LemmataRinfty} we recalled useful results on the relationship between the (non-)finiteness of Reidemeister numbers and characteristic subgroups. 
In the category of linear algebraic groups (over a fixed field $\K$), one has that $\mbU_n$ is the commutator subgroup scheme of $\mbB_n$, hence (algebraically) characteristic in $\mbB_n$. In the general abstract case one has the following.
\begin{rmk} \label{rmk:Uncharacteristic}
The subgroup $\mbU_n(R)$ is not necessarily characteristic in $\mbB_n(R)$, even over integral domains; cf.~\cite[Section~3]{Bn1}. Regardless, if $R$ is an integral domain, $\mbU_n(R)$ is the \emph{Hirsch--Plotkin radical} (hence a characteristic subgroup) of $\PB_n(R)$; see, e.g., \cite[Section~3]{Bn1} for a proof.
\end{rmk} 

\section{The theorem of Lev\v{c}uk} \label{sec:Levchuk}

The main ingredient in the proof of Theorem~\ref{thm:Additive} is a remarkable result due to Vladimir Lev\v{c}uk describing the automorphisms of the group of unitriangular matrices; see~\cite{LevchukOriginal}. 
We begin by recasting his theorem in the terminology and form needed for the present work --- in particular, we restate the results directly for \emph{upper} unitriangular matrices $\mbU_n(R)$ with the usual matrix multiplication, instead of \emph{lower} unitriangular matrices $\mbU^{-}_n(R)$ as in the paper~\cite{LevchukOriginal}. 

\begin{thm}[{Lev\v{c}uk~\cite{LevchukOriginal}}] \label{thm:Levchuk}
Let $R$ be an integral domain and let $n \geq 4$. Then any automorphism $\psi \in \Aut(\mbU_n(R))$ can be written as a product 
\[\psi = \iota \circ \mc{z} \circ \sigma \circ \tau^\veps \circ \alpha_\ast, \]
 where 
 \begin{itemize}
     \item $\iota \in \Inn(\PB_n(R))$ is conjugation by an element of $\PB_n(R) \supseteq \mbU_n(R)$; see Section~\ref{Levchuk:inner},
     \item $\mc{z} \in \mc{Z}$ is a central automorphism; see Section~\ref{Levchuk:central}, 
     \item $\sigma \in \mc{U}^{(c)}$ is an automorphism of type $\Sigma$; see Section~\ref{Levchuk:sigma}, 
     \item $\tau$ is the flip automorphism and $\veps \in \set{0,1}$; see Section~\ref{Levchuk:flip}, and 
     \item $\alpha_\ast$ is induced by a ring automorphism $\alpha \in \anel{\Aut}(R)$; see Section~\ref{Levchuk:ring}.
 \end{itemize}
\end{thm}

In what follows we elucidate the automorphisms appearing in the statement and, since the theorem does not appear in the original paper~\cite{LevchukOriginal} in the form stated above, we explain in Section~\ref{Levchuk:proof} how it follows from Lev\v{c}uk's work. 

\subsection{Conjugation by an element of $\PB_n(R) \supseteq \mbU_n(R)$}\label{Levchuk:inner}
The first elements of $\Aut(\mbU_n(R))$ in Lev\v{c}uk's list are inner automorphisms, i.e., conjugation. In a group $G$, given $g\in G$ we denote by $\iota_g \in \Aut(G)$ the inner automorphism $\iota_g(h) = ghg^{-1}$. We then have the subgroup $\Inn(\mbU_n(R)) \leq \Aut(\mbU_n(R))$ of inner automorphisms of $\mbU_n(R)$. Since $\mbU_n(R)$ is also invariant under conjugation by \emph{diagonal} matrices $d \in \mbD_n(R)$ due to the relations~\bref{rel:commutators},
it follows that the maps $\iota_d \in \Inn(\mbB_n(R))$ for $d \in \mbD_n(R)$ are also automorphisms of $\mbU_n(R)$. Thus the whole set $\Inn(\mbB_n(R))$ of inner automorphisms of $\mbB_n(R)$ is a subgroup of $\Aut(\mbU_n(R))$ in the obvious way. 

\begin{rmk} \label{obs:innerautomorphisms}
One has the decomposition $\Inn(\mbB_n(R)) = \Inn(\mbU_n(R)) \cdot \Inn(\mbD_n(R)) = \Inn(\mbD_n(R)) \cdot \Inn(\mbU_n(R))$ because $\iota_{xy} = \iota_x \circ \iota_y$ for any $x,y \in \mbB_n(R)$ and $\mbB_n(R) = \mbU_n(R) \rtimes \mbD_n(R)$. Recall moreover that $\Inn(G) \cong G/Z(G)$ for any group $G$. In particular, since $Z_n(R) \subseteq Z(\mbB_n(R))$, one has that conjugation by an equivalence class $[d] \in \mbD_n(R)/Z_n(R)$ is well-defined in the groups $\mbB_n(R)$, $\PB_n(R)$ and $\mbU_n(R)$, and coincides with conjugation by any representative $d \in [d]$. Hence 
\[\Inn(\mbB_n(R)) = \Inn(\PB_n(R)) \leq \Aut(\mbU_n(R)).\]
\end{rmk}

\subsection{Central automophisms} \label{Levchuk:central}
Next we consider automorphisms which only modify the center 
\[Z(\mbU_n(R)) = \gamma_{n-1}(\mbU_n(R)) = \mc{E}_{1,n}(R):=\langle \{ \ekl{1,n}(r) \in \mbU_n(R) \mid r \in R \}\rangle.\] 
Given an index $i \in \set{1,\ldots,n}$, any endomorphism $\lambda \in \End(\Addi(R))$ of the underlying additive group of $R$ gives rise to an element $\zeta_i(\lambda) \in \Aut(\mbU_n(R))$ as in~\cite[65]{LevchukOriginal}, defined by
\[ \zeta_i(\lambda)((a_{i,j})) = (a_{i,j}) \cdot \ekl{1,n}(\lambda(a_{i,i+1})) \, \text{ for every matrix } \, (a_{i,j}) \in \mbU_n(R). \]
The subgroup
\[ \mc{Z} = \gera{ \, \set{ \, \zeta_i(\lambda) \in \Aut(\mbU_n(R)) \, \mid \, i \in \set{1,\ldots,n}, \lambda \in \End(\Addi(R)) \, } \, } \]
is called the \emph{group of central automorphisms}.

\begin{rmk} \label{obs:centralautomorphisms}
Since the elementary matrices $\ekl{1,n}(r)$ are central in $\mbU_n(R)$, one has that central automorphisms commute with inner automorphisms of $\mbU_n(R)$, i.e., the subgroups $\mc{Z} \leq \Aut(\mbU_n(R))$ and $\Inn(\mbU_n(R)) \leq \Aut(\mbU_n(R))$ commute. (Though elements of $\mc{Z}$ need not commute with conjugation by diagonal matrices.) 
\end{rmk}

\subsection{Automorphisms of type $\Sigma$}\label{Levchuk:sigma}
We now describe what we call \emph{automorphisms of type} $\Sigma$. Suppose there exist a function $\lambda: R \to R$ and an element $a \in R$ satisfying
\begin{equation} \label{defsigmas} \lambda(r+s) = ars + \lambda(r) + \lambda(s) \, \text{ for all } \, r,s \in R. \end{equation}
Then the maps $\sigma_{\lambda,a}$ and $\sigma_{\lambda,a}'$, defined on the generators of $\mbU_n(R)$ via
\[ \sigma_{\lambda,a} : \begin{cases} e_{1,2}(r) & \mapsto e_{1,2}(r) \cdot e_{2,n}(ar) \cdot e_{1,n}(\lambda(r) - ar^2) \\ e_{i,i+1}(r) & \mapsto e_{i,i+1}(r) \, \, \text{ if } i > 1, \end{cases} \]
\[ \sigma_{\lambda,a}' : \begin{cases} e_{n-1,n}(r) & \mapsto e_{n-1,n}(r) \cdot e_{1,n-1}(ar) \cdot e_{1,n}(\lambda(r)) \\ e_{i,i+1}(r) & \mapsto e_{i,i+1}(r) \, \, \text{ if } i < n-1, \end{cases} \]
induce automorphisms of $\mbU_n(R)$ as seen in~\cite[66]{LevchukOriginal}. Abusing notation, we shall also denote such automorphisms by $\sigma_{\lambda,a}$ and $\sigma_{\lambda,a}'$. Lev\v{c}uk denotes by $\mc{U}^{(c)}$ the group of all automorphisms as above, i.e. 
\[ \mc{U}^{(c)} = \gera{ \, \set{ \,\sigma_{\lambda,a}, \sigma_{\lambda,a}' \in \Aut(\mbU_n(R)) \, \mid \, \lambda \text{ and } a \text{ satisfy Condition~\bref{defsigmas}} \, } \, }. \]
For simplicity we call $\mc{U}^{(c)} \leq \Aut(\mbU_n(R))$ the subgroup of automorphisms of type $\Sigma$, and elements of $\mc{U}^{(c)}$ will typically be denoted by $\sigma$.

\subsection{Flip automorphism}\label{Levchuk:flip}
Next in our list is the `flip' automorphism, to be denoted by $\tau \in \Aut(\mbU_n(R))$. The reader familiar with Chevalley--Demazure groups might know $\tau$ as the restriction to $\mbU_n(R)$ of the automorphism of $\SL_n(R)$ induced by the unique graph automorphism of the Dynkin diagram of type $\mathtt{A}_{n-1}$; cf. \cite[65]{LevchukOriginal} and \cite[91]{Steinberg}. In our case, $\tau$ is most conveniently described by the rule
\[ \tau(\eij(r)) = e_{n-j+1,n-i+1}((-1)^{j-i-1}r) \, \text{ for } \eij(r) \in \mbU_n(R) \]
on the generators of $\mbU_n(R)$. 
Pictorially, $\tau$ is obtained by flipping the entries of the matrices in $\mbU_n(R)$ along the anti-diagonal and adjusting the signs of the matrix entries in accordance with the commutator relations~\bref{rel:commutators}. For instance, for $n=3$ and $n=4$ the flip is define, respectively, by
\begin{align*}
\left(\begin{smallmatrix}
    1 & r & 0 \\ 0 & 1 & 0 \\ 0 & 0 & 1
\end{smallmatrix}\right)&\xmapsto{\tau}\left(\begin{smallmatrix}
    1 & 0 & 0 \\ 0 & 1 & r \\ 0 & 0 & 1
\end{smallmatrix}\right), \quad 
%
\left(\begin{smallmatrix}
    1 & 0 & 0 \\ 0 & 1 & r \\ 0 & 0 & 1
\end{smallmatrix}\right)
\xmapsto{\tau}\left(\begin{smallmatrix}
    1 & r & 0 \\ 0 & 1 & 0 \\ 0 & 0 & 1
\end{smallmatrix}\right), \quad \text{ and } 
%
\end{align*}
\begin{align*}
\left(\begin{smallmatrix}
    1 & r & 0 & 0 \\ 0 & 1 & 0 & 0 \\ 0 & 0 & 1 & 0 \\0 & 0 & 0 & 1 
\end{smallmatrix}\right)&
\xmapsto{\tau}\left(\begin{smallmatrix}
    1 & 0 & 0 & 0 \\ 0 & 1 & 0 & 0 \\ 0 & 0 & 1 & r \\ 0 & 0 & 0 & 1
\end{smallmatrix}\right), \, \,
%
\left(\begin{smallmatrix}
    1 & 0 & 0 & 0 \\ 0 & 1 & r & 0 \\ 0 & 0 & 1 & 0 \\0 & 0 & 0 & 1 
\end{smallmatrix}\right)
\xmapsto{\tau}\left(\begin{smallmatrix}
    1 & 0 & 0 & 0 \\ 0 & 1 & r & 0 \\ 0 & 0 & 1 & 0 \\ 0 & 0 & 0 & 1
\end{smallmatrix}\right), \, \,
\left(\begin{smallmatrix}
    1 & 0 & 0 & 0 \\ 0 & 1 & 0 & 0 \\ 0 & 0 & 1 & r \\0 & 0 & 0 & 1 
\end{smallmatrix}\right)
\xmapsto{\tau}\left(\begin{smallmatrix}
    1 & r & 0 & 0 \\ 0 & 1 & 0 & 0 \\ 0 & 0 & 1 & 0 \\ 0 & 0 & 0 & 1
\end{smallmatrix}\right). 
\end{align*}
It is immediate that $\tau^2 = \id$.

\subsection{Group automorphisms induced by ring automorphisms}\label{Levchuk:ring}
To finish our list of automorphisms of $\mbU_n(R)$ recall that, since $\mbU_n$ is an affine $\Z$-subscheme of $\GL_n$, any ring automorphism $\alpha \in \anel{\Aut}(R)$ induces by functoriality a group automorphism $\alpha_\ast : \mbU_n(R) \to \mbU_n(R)$. More explicitly, given $\alpha \in \anel{\Aut}(R)$ and a matrix $(a_{ij}) \in \mbU_n(R)$, one has $\alpha_\ast((a_{ij})) := (\alpha(a_{ij}))$. With this in mind we consider $\anel{\Aut}(R)$ as a subgroup of $\Aut(\mbU_n(R))$ in the obvious way and call it the \emph{subgroup of ring automorphisms} of $\mbU_n(R)$.

\begin{rmk} \label{obs:ringautomorphisms}
All subgroups $\sbgpeij(R) \leq \mbU_n(R)$ are $\alpha_\ast$-invariant for any ring automorphism $\alpha \in \anel{\Aut}(R)$. This is because ring automorphisms preserve~$1$ and so $\alpha_\ast(\eij(r)) = \eij(\alpha(r))$. In particular, $\alpha_\ast$ induces an automorphism $\barra{\alpha}_\ast$ on the abelianization $\mbU_n(R)^{\mathrm{ab}}$ such that the image of every $\mc{E}_{i,i+1}(R)$ in $\mbU_n(R)^{\mathrm{ab}}$ is also $\barra{\alpha}_\ast$-invariant. Moreover, since $\sbgpeij(R) \cong \Addi(R)$ we may write $\alpha_\ast(\eij(r)) = \eij(\addi{\alpha}(r))$, where $\addi{\alpha}$ is the same map $\alpha$ but viewed as an automorphism of the underlying additive group $\Addi(R) = (R,+)$. In particular, the restriction $\alpha_\ast|_{\sbgpeij(R)}$ is also interpreted as an automorphism of $\Addi(R)$.
\end{rmk}

\subsection{Proof of Theorem~\ref{thm:Levchuk}}\label{Levchuk:proof}

Here we spell out how Theorem~\ref{thm:Levchuk} follows from Lev\v{c}uk's work. We want to show that $\Aut(\mbU_n(R))$ can be written as a product
\[\Aut(\mbU_n(R)) = \Inn(\PB_n(R)) \cdot \mc{Z} \cdot \mc{U}^{(c)} \cdot \gera{\tau} \cdot \anel{\Aut}(R)\]
whenever $n\geq 4$, which implies that any $\psi \in \Aut(\mbU_n(R))$ can be written as a product 
\[\psi = \iota \circ \mc{z} \circ \sigma \circ \tau^\veps \circ \alpha_\ast,\] with $\iota \in \Inn(\PB_n(R))$, $\mc{z} \in \mc{Z}$, $\sigma \in \mc{U}^{(c)}$, $\varepsilon \in \{0,1\}$, $\tau$ is the flip automorphism and $\alpha_\ast$ is induced by a ring automorphism. We need to address the cases $n\geq 5$ and $n=4$ separately.

Write $D \leq \Aut(\mbU_n(R))$ for the subgroup of automorphisms which are given by conjugation by a diagonal matrix. By~\cite[Corollary~3]{LevchukOriginal}, if $n \geq 5$ then there exist two further subgroups $W$ and $V$ of $\Aut(\mbU_n(R))$ such that $\Aut(\mbU_n(R))$ decomposes as 
\[ \Aut(\mbU_n(R)) = (((\mc{Z} \cdot \Inn(\mbU_n(R)) \cdot \mc{U}^{(c)} \cdot V) \ltimes W) \ltimes D) \ltimes \anel{\Aut}(R). \]
(In \cite{LevchukOriginal}, the group $\Inn(\mbU_n(R))$ is denoted by a calligraphic J.)

The generators of the subgroup $V \leq \Aut(\mbU_n(R))$ above are induced by the assignments $\eta_a$ defined in~\cite[66]{LevchukOriginal}, which in turn are defined by choosing elements $a \in R$ belonging ([\emph{loc. cit.}]) to the annihilator of the set 
\[\set{r\cdot s \mid r,s \in R} \cup \set{2} \subseteq R.\] 
(Notice that we can ignore the maps of the form $\eta_b$ from~\cite[66]{LevchukOriginal} because our ring $R$ is commutative.) Moreover one readily checks that $\eta_a$ induces the identity map if $a = 0$. Since our ring $R$ is an integral domain, the only possible choice for annihilator in the above is $a = 0$, hence $V = \set{\id}$.

The generators of the subgroup $W \leq \Aut(\mbU_n(R))$ are the so-called idempotent automorphisms, described in~\cite[65]{LevchukOriginal}. The definition of such a generator $\tau_e \in W$ depends on the choice of an idempotent element $e$ lying in the center $Z(R)$ of the ring $R$ [\emph{loc. cit.}]. Here we stress that the flip automorphism $\tau$ coincides with the generator $\tau_0 \in W$ associated to the trivial idempotent $e = 0$. Since $R$ is an integral domain we have $Z(R) = R$ and it furthermore only admits the trivial idempotents $0$ and $1$. However, the description in~[\emph{loc. cit.}] also shows that $\tau_1$ induces the identity map. Thus in our case $W = \gera{\tau}$, which is a group of order $2$.

Therefore \cite[Corollary~3]{LevchukOriginal} applied to integral domains actually yields
\begin{equation} \label{eq:Levchuk} \Aut(\mbU_n(R)) = (((\mc{Z} \cdot \Inn(\mbU_n(R)) \cdot \mc{U}^{(c)}) \ltimes \gera{\tau}) \ltimes D) \ltimes \anel{\Aut}(R) \end{equation}
in the case $n \geq 5$.

Now, recalling that a semi-direct product $G = N \rtimes Q$ can be equivalently written as $G = Q \ltimes N$ and moreover that $G = N \cdot Q = Q \cdot N$, we may rewrite equality~\bref{eq:Levchuk} in the desired way. Indeed, since $\mc{Z}$ and $\Inn(\mbU_n(R))$ commute by Remark~\ref{obs:centralautomorphisms}, equality~\bref{eq:Levchuk} gives 
\[ \Aut(\mbU_n(R)) = ( D \rtimes ((\Inn(\mbU_n(R)) \cdot \mc{Z} \cdot \mc{U}^{(c)}) \ltimes \gera{\tau}) ) \ltimes \anel{\Aut}(R). \]
As in Remark~\ref{obs:innerautomorphisms}, the group $\Inn(\mbB_n(R)) = \Inn(\PB_n(R))$ is viewed as a subgroup of $\Aut(\mbU_n(R))$ in the obvious way. Moreover it is the product of $\Inn(\mbU_n(R))$ with $\Inn(\mbD_n(R) / Z_n(R))$, the latter being viewed as the set of automorphisms of $\mbU_n(R)$ which are conjugation by a diagonal matrix modulo multiples of the identity (which are central in $\mbB_n(R)$). One thus has $D = \Inn(\mbD_n(R) / Z_n(R))$ and $\Inn(\PB_n(R)) = \Inn(\mbU_n(R)) \cdot D = D \cdot \Inn(\mbU_n(R))$ as a subgroup of $\Aut(\mbU_n(R))$. In particular, equality~\bref{eq:Levchuk} becomes 
\[\Aut(\mbU_n(R)) = \Inn(\PB_n(R)) \cdot \mc{Z} \cdot \mc{U}^{(c)} \cdot \gera{\tau} \cdot \anel{\Aut}(R),\]
as desired.

We now address the case $n = 4$. The starting point is a description of $\Aut(\mbU_4(R))$ which is slightly similar to the one from the case $n \geq 5$. 
More precisely, by \cite[Theorem~2]{LevchukOriginal} there exists a subgroup $\widetilde{S} \leq \Aut(\mbU_4(R))$ such that $\Aut(\mbU_4(R))$ decomposes as 
\[ \Aut(\mbU_4(R)) = ((\mc{Z} \cdot \Inn(\mbU_4(R)) \cdot \mc{U}^{(c)}) \ltimes (\widetilde{S} \cdot D)) \ltimes \anel{\Aut}(R).\] 
We take a closer look at how the subgroup $\widetilde{S} \leq \Aut(\mbU_4(R))$ is defined. 
In~\cite[pp.~73 and~74]{LevchukOriginal}, Lev\v{c}uk observes that a matrix $s = \left( \begin{smallmatrix} a_{1,1} & a_{1,2} \\ a_{2,1} & a_{2,2} \end{smallmatrix} \right) \in \SL_2(R)$ gives rise to an automorphism $\widetilde{s} \in \Aut(\mbU_4(R))$ as long as it satisfies the conditions 
\[ 2a_{1,1}a_{1,2} = 2a_{2,1}a_{2,2} = 0 \, \, \text{ and } \, \, a_{i,1}a_{i,2}(x^2-x)(y^2-y)=0 \]
for $i\in\{1,2\}$ and $x,y\in R$. Since $R$ is an integral domain, the first condition implies that any such matrix $s = (a_{i,j}) \in \SL_2(R)$ must have the form
\begin{equation} \label{eq:s-Automorphisms}
s = \left( \begin{smallmatrix} a_{1,1} & 0 \\ 0 & a_{2,2} \end{smallmatrix} \right) = \left( \begin{smallmatrix} a_{1,1} & 0 \\ 0 & a_{1,1}^{-1} \end{smallmatrix} \right)  \, \text{ or } \, s = \left( \begin{smallmatrix} 0 & a_{1,2} \\ -a_{2,1} & 0 \end{smallmatrix} \right) = \left( \begin{smallmatrix} 0 & a_{1,2} \\ -a_{1,2}^{-1} & 0 \end{smallmatrix} \right).
\end{equation}
In particular, a matrix as in condition~\bref{eq:s-Automorphisms} automatically satisfies the second condition above for any $i$ and any $x,y \in R$. Going further, the automorphism $\widetilde{s} \in \Aut(\mbU_4(R))$ induced by the matrix $s = (a_{i,j})$ acts on the generators $e_{k,\ell}(r) \in \mbU_4(R)$ by the following rules. 
\[
\widetilde{s} : 
\begin{cases} e_{1,2}(r) & \mapsto e_{1,2}(a_{1,1}r) \cdot e_{3,4}(a_{1,2}r), \\ 
e_{2,3}(r) & \mapsto e_{2,3}(r), \\ 
e_{3,4}(r) & \mapsto e_{1,2}(a_{2,1}r) \cdot e_{2,2}(a_{2,2}r), \\
e_{1,3}(r) & \mapsto e_{1,3}(a_{1,1}r) \cdot e_{2,4}(-a_{1,2}r), \\
e_{2,4}(r) & \mapsto e_{1,3}(-a_{2,1}r) \cdot e_{2,4}(a_{2,2}r), \\
e_{1,4}(r) & \mapsto e_{1,4}((a_{1,1}a_{2,2}+a_{1,2}a_{2,1})r).
\end{cases}
\]
The subgroup $\widetilde{S} \leq \Aut(\mbU_4(R))$ is precisely the group generated by all automorphisms $\widetilde{s}$ as above~[\emph{loc. cit.}]. 

The point now is that the product $\widetilde{S} \cdot D$ is isomorphic to $\gera{\tau} \ltimes D$, where $\tau$ is the flip automorphism. In effect, since a matrix $s$ giving rise to $\widetilde{s} \in \widetilde{S}$ must be as in condition~\bref{eq:s-Automorphisms}, the action of $\widetilde{s}$ given previously reduces as follows.
\[
\text{For } s=\left( \begin{smallmatrix} a_{1,1} & 0 \\ 0 & a_{1,1}^{-1} \end{smallmatrix} \right) \text{ one has } \, \widetilde{s}: \begin{cases} e_{1,j}(r) & \mapsto e_{1,j}(a_{1,1}r) \, \text{ if } j \in \{ 2,3 \}, \\ 
e_{k,\ell}(r) & \mapsto e_{k,\ell}(r) \, \text{ if } (k,\ell) \in \{ (2,3), (1,4)\}, \\ 
e_{i,4}(r) & \mapsto e_{i,4}(a_{1,1}^{-1}r) \, \text{ if } i \in \{ 2,3 \},
\end{cases}
\]
\[
\text{whereas for } \, s=\left( \begin{smallmatrix} 0 & a_{1,2} \\ -a_{1,2}^{-1} & 0 \end{smallmatrix} \right) \, \text{ one has } \, \widetilde{s} : 
\begin{cases} e_{1,2}(r) & \mapsto e_{3,4}(a_{1,2}r), \\ 
e_{2,3}(r) & \mapsto e_{2,3}(r), \\ 
e_{3,4}(r) & \mapsto e_{1,2}(-a_{1,2}^{-1}r), \\
e_{1,3}(r) & \mapsto e_{2,4}(-a_{1,2}r), \\
e_{2,4}(r) & \mapsto e_{1,3}(a_{1,2}^{-1}r), \\
e_{1,4}(r) & \mapsto e_{1,4}(-r).
\end{cases}
\]
Using the above descriptions, direct matrix computations show that the map $\widetilde{s}$ is in fact a composition of the flip automorphism and conjugation by a diagonal matrix. More precisely, it holds 
\[\widetilde{s} = \tau \circ \iota_d,\]
where $\iota_d \in D$ is conjugation by the diagonal matrix 
\[ d = 
\begin{cases} \left( \begin{smallmatrix} a_{1,1} & & & \\ & 1 & & \\ & & 1 & \\ & & & a_{1,1} \end{smallmatrix} \right), & \text{ in case } s = \left( \begin{smallmatrix} a_{1,1} & 0 \\ 0 & a_{1,1}^{-1} \end{smallmatrix} \right), \\
\left( \begin{smallmatrix} a_{1,2} & & & \\ & 1 & & \\ & & 1 & \\ & & & -a_{1,2} \end{smallmatrix} \right), & \text{ in case } s = \left( \begin{smallmatrix} 0 & a_{1,2} \\ -a_{1,2}^{-1} & 0 \end{smallmatrix} \right).
\end{cases}
\]
Thus the only generator of $\widetilde{S} \cdot D$ not lying in $D$ is the flip automorphism. One readily checks that $D$ is normal in $\widetilde{S} \cdot D$, and that the product splits as $\gera{\tau} \ltimes D$. 

Hence, similarly to the case $n \geq 5$, one obtains the same equality~\bref{eq:Levchuk}
\[ \Aut(\mbU_4(R)) = (((\mc{Z} \cdot \Inn(\mbU_4(R)) \cdot \mc{U}^{(c)}) \ltimes \gera{\tau}) \ltimes D) \ltimes \anel{\Aut}(R). \]
The final steps of the case $n\geq 5$ thus also carry over when $n = 4$, yielding 
\[\Aut(\mbU_4(R)) = \Inn(\PB_4(R)) \cdot \mc{Z} \cdot \mc{U}^{(c)} \cdot \gera{\tau} \cdot \anel{\Aut}(R),\]
which finishes the proof. \qed

\section{Restatement and proof of Theorem~\ref{thm:Additive}}\label{sec:provadothmB} 
Before recalling our main technical theorem we need a bit more notation. Given a \emph{ring} automorphism $\alpha \in \anel{\Aut}(R)$ we consider the following automorphisms on the underlying additive group $\Addi(R) = (R,+)$. 
Firstly, we let $\addi{\alpha} \in \Aut(\Addi(R))$ denote the same automorphism $\alpha$ now viewed as a \emph{group} automorphism of $\Addi(R)$, i.e. $\addi{\alpha}(r)$ is just $\alpha(r)$ for any $r \in R$. Secondly, in the direct product $\Addi(R) \times \Addi(R)$ of two copies of $(R,+)$ we let $\tau_\alpha \in \Aut(\Addi(R) \times \Addi(R))$ denote the `flip' automorphism induced by~$\alpha$, that is, $\tau_\alpha((r,s)) = (\addi{\alpha}(s),\addi{\alpha}(r))$.

\theoremstyle{plain}
\newtheorem*{thmB}{Theorem~\ref{thm:Additive}}
\begin{thmB} Let $R$ be an integral domain with a \emph{finitely generated} group of units $\Mult(R)$. Assume that both $R(\addi{\alpha})$ and $R(\tau_\alpha)$ are infinite for all $\alpha \in \anel{\Aut}(R)$, where $\addi{\alpha}$ and $\tau_\alpha$ are as above. Then the groups $\mbB_n(R)$ and $\PB_n(R)$ have property $\Ri$ for all $n \geq 4$.
\end{thmB}

For the remaining of Section~\ref{sec:provadothmB}, we assume that $R$ is as in the statement and $n \geq 4$. 

Because $\PB_n(R)$ is a characteristic quotient of $\mbB_n(R)$, Lemma~\ref{lem:desempre} shows that $\mbB_n(R)$ will have property $\Ri$ if $\PB_n(R)$ does so. To prove Theorem~\ref{thm:Additive} it thus suffices to show that all $\phee \in \Aut(\PB_n(R))$ satisfy $R(\phee)=\infty$. This analysis will be done in several steps. The first observation is the following. 

\begin{lem}\label{lem:useless}
Let $G$ be a group and let $H \leq G$ be a characteristic subgroup. Suppose there exists a subgroup $A \leq \Aut(H)$ such that $\Aut(H) = \Inn(G) \cdot A$, where $\Inn(G)$ is viewed as a subgroup of $\Aut(H)$ in the obvious way. Then any automorphism $\phee \in \Aut(G)$ can be written as a product $\phee = \iota \circ \psi$ where $\iota \in \Inn(G)$ and $\psi|_H \in A$.
\end{lem}

\begin{proof}
Because $H$ is characteristic, it holds $\Inn(H) \leq \Inn(G) \leq \Aut(H)$ in the obvious way and, given any automorphisms $\Phi, \Psi \in \Aut(G)$, one has $(\Phi \circ \Psi)|_H = \Phi|_H \circ \Psi|_H$. 

Let $\phee \in \Aut(G)$. Pick a right transversal $T \subseteq \Aut(G)$ for $\Inn(G)$, so that $\phee = \iota_0 \circ \tau$ for some $\iota_0 \in \Inn(G)$ and $\tau \in T$. By hypothesis, the restriction $\tau|_H$ can be written as $\tau|_H = \iota_1 \circ \alpha$ where $\iota_1 \in \Inn(G)$ and $\alpha \in A$. Define $\iota = \iota_0 \circ \iota_1 \in \Inn(G)$ and $\psi = \iota_1^{-1} \circ \tau \in \Aut(G)$. One then has
\[ \phee = \iota_0 \circ \tau = (\iota_0 \circ \iota_1) \circ (\iota_1^{-1} \circ \tau) = \iota \circ \psi \]
and $\psi|_H = \iota_1^{-1}|_H \circ \tau|_H = \iota_1^{-1} \circ \iota_1 \circ \alpha = \alpha \in A$, as desired.
\end{proof}

Lemma~\ref{lem:useless} thus assures that we can write each $\phee \in \Aut(\PB_n(R))$ as $\phee = \iota \circ \psi$ with $\iota \in \Inn(\PB_n(R))$ and where the restriction $\psi'$ of $\psi \in \Aut(\PB_n(R))$ to the characteristic subgroup $\mbU_n(R)$ is described by Theorem~\ref{thm:Levchuk}. The next steps, given in Section~\ref{sec:gpA}, shall yield $R(\psi')=\infty$. Finally, we use this equality in Section~\ref{sec:provaThmB} to show that $R(\phee) = \infty$. 

\subsection{A subset of automorphisms of $\mbU_n(R)$}\label{sec:gpA} 
As shown in Theorem~\ref{thm:Levchuk}, one may write
\[\Aut(\mbU_n(R)) = \Inn(\PB_n(R)) \cdot \mc{Z} \cdot \mc{U}^{(c)} \cdot \gera{\tau} \cdot \anel{\Aut}(R). \] 
In this section, we consider the subgroup $\mc{A} = \mc{Z} \cdot \mc{U}^{(c)} \cdot \gera{\tau} \cdot \anel{\Aut}(R) \leq \Aut(\mbU_n(R))$. We claim that, under the hypothesis of Theorem~\ref{thm:Additive}, 
\begin{align} \label{eq:psilinha}
R(\psi') = \infty \, \text{ for all } \, \psi' \in \mc{A}.
\end{align}

By definition, an automorphism $\psi' \in \mc{A}$, is of the form $ \psi' = \mc{z} \circ \sigma \circ \tau^\veps \circ \alpha_\ast$ with  $\veps \in \set{0,1}$, where $\mc{z} \in \mc{Z}$, $\sigma \in \mc{U}^{(c)}$, $\alpha_\ast \in \anel{\Aut}(R)$; cf. Section~\ref{sec:Levchuk}.

Since $\gamma_2(\mbU_n(R)) = [\mbU_n(R), \mbU_n(R)]$ is characteristic in $\mbU_n(R)$, the map $\psi'$ induces an automorphism $\barra{\psi'} = \barra{\mc{z} \circ \sigma \circ \tau^\veps \circ \alpha_\ast}$ on the abelianization $\mbU_n(R)^{\mathrm{ab}}$ and moreover $R(\psi') \geq R(\barra{\psi'})$; c.f.\ Lemma~\ref{lem:desempre}. It suffices to check $R(\barra{\psi'}) = \infty$.
The following two lemmata simplify this task by assuring that $\barra{\mc{z} \circ \sigma} = \id$, whence $\barra{\psi'} = \barra{\tau^\veps \circ \alpha_\ast}$ and $R(\barra{\psi'}) = R(\barra{\tau^\veps \circ \alpha_\ast})$.

\begin{lem} \label{obs:centralautomorphisms2}
For $n \geq 3$, any central automorphism $\mc{z} \in \mc{Z} \leq \Aut(\mbU_n(R))$ induces the identity on the abelianization $\mbU_n(R)^{\mathrm{ab}}$.
\end{lem}

\begin{proof}
In fact a stronger statement holds: $\mc{z}$ induces the identity on each factor $\gamma_k(\mbU_n(R)) / \gamma_{k+1}(\mbU_n(R))$ for $k < n-1$. This is immediate from the definition since $\mc{z}$ acts as multiplication by an element lying in the center $Z(\mbU_n(R)) = \mc{E}_{1,n}(R)$. (Recall that $\mc{z}$ induces an automorphism $\barra{\mc{z}}_k$ of the quotient $\gamma_k(\mbU_n(R)) / \gamma_{k+1}(\mbU_n(R))$ because the lower central series $\gamma_i(\mbU_n(R))$ is a characteristic series.) In particular, $\barra{\mc{z}}_1 = \id$ on the abelianization $\mbU_n(R)^{\mathrm{ab}} = \gamma_1(\mbU_n(R)) / \gamma_{2}(\mbU_n(R)) = \mbU_n(R) / [\mbU_n(R),\mbU_n(R)]$.
\end{proof}

\begin{lem} \label{obs:sigmaautomorphisms}
For $n \geq 3$, any automorphism of type $\Sigma$ induces the identity on the abelianization $\mbU_n(R)^{\mathrm{ab}}$.
\end{lem}

\begin{proof}
This is also straightforward from the definition, for $e_{2,n}(r)$, $e_{1,n-1}(r)$ and $e_{1,n}(r)$ all belong to $\gamma_2(\mbU_n(R)) = [\mbU_n(R),\mbU_n(R)]$ for all $r \in R$, whence any $\sigma \in \mc{U}^{(c)}$ induces a trivial action on $\mbU_n(R)^{\mathrm{ab}} = \mbU_n(R) / [\mbU_n(R), \mbU_n(R)] = \gamma_1(\mbU_n(R)) / \gamma_2(\mbU_n(R))$. 
\end{proof}

In the following step --- Proposition~\ref{obs:flip} --- we construct a $\barra{\tau^\veps \circ \alpha_\ast}$-invariant subgroup $N$ of the abelianization $\mbU_n(R)^{\mathrm{ab}}$ such that the quotient $Q=\mbU_n(R)^{\mathrm{ab}}/N$ is also $\barra{\tau^\veps \circ \alpha_\ast}$-invariant. Then, in Proposition~\ref{pps:autA}, we derive equality~\bref{eq:psilinha} by showing that the induced automorphism $\barra{\barra{\tau^\veps \circ \alpha_\ast}}$ on $Q$ is such that $R(\barra{\barra{\tau^\veps \circ \alpha_\ast}})=\infty$, hence $R(\barra{\tau^\veps \circ \alpha_\ast}) 
=\infty$ by Lemma~\ref{lem:desempre}.


\begin{pps}\label{obs:flip} The abelianization $\mbU_n(R)^{\mathrm{ab}}$ has two subgroups --- $E_{\mathrm{mid}}$ and $C_{\mathrm{mid}}$ --- that are 
{$\barra{\tau^\veps \circ \alpha_\ast}$-invariant}. In particular, the map 
{$\barra{\tau^\veps \circ \alpha_\ast}$} further induces automorphisms on both such groups.
\end{pps}

\begin{proof}
We start by constructing the `middle' subgroup $E_{\mathrm{mid}}$. By definition, an element of $\prod_{i=1}^{n-1} \mc{E}_{i,i+1}(R)$ is of the form $\prod_{i=1}^{n-1}e_{i,i+1}(r_{i,i+1})$. For simplicity, we write 
\[\mathbf{e}_m(r) := \prod_{i=1}^{n-1}e_{i,i+1}(r_{i,i+1}),\] whenever $r_{i,i+1}=0$ for all $i \neq m$ and $r_{m,m+1}=r$.

We define the `middle' subgroup by setting 
\[ E_{\mathrm{mid}} = \langle\{ \mathbf{e}_m(r), \mathbf{e}_{n - m}(r) \mid r \in R \} \rangle \leq \prod_{i=1}^{n-1} \mc{E}_{i,i+1}(R),\]
where $m = \left\lceil \frac{n-1}{2} \right\rceil$. 
That is, $E_{\mathrm{mid}}$ is just the image of $\mc{E}_{m, m + 1}(R) \cdot \mc{E}_{n - m, n - m + 1}(R)$ in $\mbU_n(R)^{\mathrm{ab}}$.

For example --- and to justify the notation `middle' --- consider the cases $n = 4$ and $n = 5$. In the former, one has $m = \left\lceil \frac{3}{2} \right\rceil  = 2$ but also $n - m = 2$, so that $\mc{E}_{m, m + 1}(R) = \mc{E}_{n - m, n - m + 1}(R) = \mc{E}_{2, 3}(R)$, hence $E_{\mathrm{mid}}$ is isomorphic to the elementary subgroup $\mc{E}_{2,3}(R)$ found in the the middle of the superdiagonal of $\mbU_4(R)$. In case $n = 5$, one has $m = \left\lceil \frac{4}{2} \right\rceil  = 2$ and $n-m=5-2=3$, so that $E_{\mathrm{mid}}$ is the direct product $\mc{E}_{2,3}(R) \times \mc{E}_{3,4}(R)$ of the two `middle' subgroups of the superdiagonal of $\mbU_5(R)$.

The middle subgroup $E_{\mathrm{mid}}$ has the group $C_{\mathrm{mid}} = \prod_{i=1}^{n-1} E_{i,i+1}$ as (group-theoretic) complement, where 
\[ E_{i,i+1} = \begin{cases} \mc{E}_{i,i+1}(R), & \text{ if } i \notin \set{ m, \, n - m }, \\
\set{1} & \text{ otherwise}. \end{cases} \]

Let us have a closer look at the automorphism $\barra{\tau}$ induced by the flip automorphism $\tau$ on the abelianization $\mbU_n(R)^{\mathrm{ab}} \cong \prod_{i=1}^{n-1} \mc{E}_{i,i+1}(R)$.
We observe that $\tau$ does not alter the signs of the first superdiagonal entries of matrices of $\mbU_n(R)$. It also maps any subgroup $\mc{E}_{i,i+1}(R)$ on the superdiagonal isomorphically onto another subgroup $\mc{E}_{j,j+1}(R)$ on the superdiagonal.

Using these facts, one concludes that $\barra{\tau}$ acts on $\mbU_n(R)^{\mathrm{ab}}$ via 
\begin{align} \label{tau1} \begin{split} \barra{\tau}( (e_{1,2}(r_1), \, e_{2,3}(r_2), \ldots, \, e_{n-2,n-1}(r_{n-2}), \, e_{n-1,n}(r_{n-1}) )) & = \\
 (e_{1,2}(r_{n-1}), \, e_{2,3}(r_{n-2}), \ldots, \, e_{n-2,n-1}(r_{2}), \, e_{n-1,n}(r_{1}) ) & \end{split} \end{align}
 flipping coordinates, with $\mbU_n(R)^{\mathrm{ab}}$ written as the direct product $\mc{E}_{1,2}(R) \times \mc{E}_{2,3}(R) \times \cdots \times \mc{E}_{n-1,n}(R)$.


The action~\bref{tau1} of the flip automorphism $\tau$ thus shows that both $E_{\mathrm{mid}}$ and $C_{\mathrm{mid}}$ are $\barra{\tau}$-invariant. The fact that both $E_{\mathrm{mid}}$ and $C_{\mathrm{mid}}$ are $\overline{\alpha_\ast}$-invariant follows from Remark~\ref{obs:ringautomorphisms}.
\end{proof}

\begin{pps}\label{pps:autA} Under the hypothesis of Theorem~\ref{thm:Additive}, each automorphism $\psi' \in \mc{A}$ satisfies $R(\psi')=\infty$. 
\end{pps}
\begin{proof}
Write $ \psi' = \mc{z} \circ \sigma \circ \tau^\veps \circ \alpha_\ast \in \mc{A}$ with  $\veps \in \set{0,1}$, where $\mc{z} \in \mc{Z}$, $\sigma \in \mc{U}^{(c)}$, $\alpha_\ast \in \anel{\Aut}(R)$.  
 As shown in Lemmata~\ref{obs:centralautomorphisms2} and~\ref{obs:sigmaautomorphisms} it holds $\barra{\mc{z} \circ \sigma} = \id$, whence $\barra{\psi'} = \barra{\tau^\veps \circ \alpha_\ast}$.

Following (the proof of) Proposition~\ref{obs:flip}, consider the `middle' subgroup $E_{\mathrm{mid}} \leq \mbU_n(R)^{\mathrm{ab}}$, i.e., the image in $\mbU_n(R)^{\mathrm{ab}}$ of the product 
\[\mc{E}_{m, m + 1}(R) \cdot \mc{E}_{n-m, n - m + 1}(R) \leq \mbU_n(R),\] 
where $m = \left\lceil \frac{n-1}{2} \right\rceil$. 
Since its complement $C_{\mathrm{mid}}$ in $\mbU_n(R)^{\mathrm{ab}}$ is also $\barra{\tau^\veps \circ \alpha_\ast}$-invariant, modding it out yields an induced automorphism $\barra{\barra{\tau^\veps \circ \alpha_\ast}}$ on $E_{\mathrm{mid}}$ and, as usual, $R(\barra{\tau^\veps \circ \alpha_\ast}) \geq R\left( \barra{\barra{\tau^\veps \circ \alpha_\ast}} \right)$ by Lemma~\ref{lem:desempre}. 

Now suppose $n$ is even. Then $m=\left\lceil \frac{n-1}{2} \right\rceil = n - \left\lceil \frac{n-1}{2} \right\rceil=n-m$, in which case the middle subgroup $E_{\mathrm{mid}}$ consists of a single copy $\mc{E}_{m, m + 1}(R)$. That is, 
\[ E_{\mathrm{mid}} = \mc{E}_{m, m + 1}(R) \cong \Addi(R), \] 
and moreover $\barra{\barra{\tau^\veps}} = \id$. In case $n$ is odd, one has $n-m=n-\left\lceil \frac{n-1}{2} \right\rceil = n - \frac{n-1}{2} = \frac{n-1}{2}+1 = \left\lceil \frac{n-1}{2} \right\rceil + 1=m+1$. Thus, $E_{\mathrm{mid}}$ is the direct product
\[ E_{\mathrm{mid}} \cong \mc{E}_{m, m + 1}(R) \times \mc{E}_{m + 1, m + 2}(R) \cong \Addi(R) \times \Addi(R). \]
Using the identification above, if $\veps = 1$ it follows from the action~\bref{tau1} shown in Section~\ref{Levchuk:flip} that $\barra{\barra{\tau^\veps}} = \barra{\barra{\tau}}$ acts on $E_{\mathrm{mid}} \cong \Addi(R) \times \Addi(R)$ by the coordinate flip 
\[
\xymatrix@R=2mm{
\barra{\barra{\tau^\veps}} : \Addi(R) \times \Addi(R) \ar[r] & \Addi(R) \times \Addi(R) \\
(r,s) \ar@{|->}[r] & (s,r).
}
\]

Writing $\alpha \in \anel{\Aut}(R)$ for the ring automorphism which induces $\alpha_\ast$, it follows from the above discussion that the induced automorphism $\barra{\barra{\tau^\veps \circ \alpha_\ast}}$ acts on the middle subgroup $E_{\mathrm{mid}}$ as follows.
\[ \barra{\barra{\tau^\veps \circ \alpha_\ast}}(r) = \addi{\alpha}(r) \, \text{ if } E_{\mathrm{mid}} \cong \Addi(R) \, \text{ (case } n \text{ even), otherwise} \]
\[ \barra{\barra{\tau^\veps \circ \alpha_\ast}}(r,s) = \begin{cases} (\addi{\alpha}(r), \addi{\alpha}(s)) & \text{ if } \veps = 0, \\ (\addi{\alpha}(s), \addi{\alpha}(r)) & \text{ if } \veps=1, \end{cases} \text{ if } E_{\mathrm{mid}} \cong \Addi(R) \times \Addi(R).  \]
In the first two cases above one has $R(\barra{\barra{\tau^\veps \circ \alpha_\ast}}) \geq R(\addi{\alpha})$, whereas in the third case ($n$ odd and $\veps=1$) it even holds $\barra{\barra{\tau^\veps \circ \alpha_\ast}} = \tau_\alpha$, where $\addi{\alpha}$ and $\tau_\alpha$ are the automorphisms defined for the statement of Theorem~\ref{thm:Additive}. Altogether, the hypotheses of the theorem yield $R\left(\barra{\barra{\tau^\veps \circ \alpha_\ast}}\right) = \infty$, whence the equality $R(\barra{\psi'}) = \infty$, as desired.    
\end{proof}

\subsection{Concluding the proof of Theorem~\ref{thm:Additive}}\label{sec:provaThmB}

As explained, it suffices to show that $\PB_n(R)$ has property~$\Ri$. Let $\phee \in \Aut(\PB_n(R))$.

Suppose first that $R^\times=\set{1}$. Then $\PB_n(R) = \mbU_n(R)$ by definition. Moreover the subgroup $D \leq \Aut(\mbU_n(R))$ of automorphisms given by conjugation by a diagonal matrix is also trivial. It follows from Theorem~\ref{thm:Levchuk} that an arbitrary automorphism $\phee \in \Aut(\PB_n(R)) = \Aut(\mbU_n(R))$ is of the form $\phee = \iota \circ \psi$ where $\iota \in \Inn(\PB_n(R)) = \Inn(\mbU_n(R))$ and $\psi \in \mc{A} \leq \Aut(\mbU_n(R))$. Thus 
{Proposition~\ref{pps:autA}} and Lemma~\ref{lem:ignoreinner} give $R(\phee)=R(\psi) = \infty$. 

Assume now that the integral domain $R$ has at least two units. 
Apply Lemma~\ref{lem:useless} taking $G = \PB_n(R)$, $H = \mbU_n(R)$ and $A = \mc{A}$ as in Section~\ref{sec:gpA}. Again by Lemma~\ref{lem:ignoreinner} we have $R(\phee) = R(\iota \circ \psi) = R(\psi)$. It thus suffices to show that $R(\psi) = \infty$.

Recall that $\psi \in \Aut(\PB_n(R))$ induces an automorphism $\barra{\psi} \in \Aut\left(\frac{\mbD_n(R)}{Z_n(R)}\right)$ on the diagonal and $\psi' := \psi|_{\mbU_n(R)} \in \mc{A} \leq \Aut(\mbU_n(R))$. If $R(\barra{\psi}) = \infty$, Lemma~\ref{lem:desempre} already yields $R(\psi) = \infty$. Assume otherwise that $R(\barra{\psi}) < \infty$. Because $\Mult(R)$ and thus $\frac{\mbD_n(R)}{Z_n(R)}$ is finitely generated, we have that $\barra{\psi}$ has finitely many fixed points by Lemma~\ref{lem:ReidemeisterAbelian}. Since $R(\psi') = \infty$ by 
{Proposition~\ref{pps:autA}} it follows from Lemma~\ref{lem:desempre} that $R(\psi) = \infty$ as well, which finishes off the proof. \qed

\begin{rmk} \label{teoremadoTimur}
The reader familiar with results on property $\Ri$ for nilpotent groups might recall a similar theorem due to T.~Nasybullov, namely~\cite[Theorem~3]{TimurUniTri}. It is stated that, if $R$ is an infinite integral domain which is \emph{finitely generated} as a $\Z$-module, then the unipotent subgroup $\mbU_n(R)$ has property $\Ri$ whenever $n > 2|R^\times|$. In particular, combined with Lemma~\ref{lem:desempre}, Nasybullov's theorem implies some special cases of Theorem~\ref{thm:Additive}. 

We point out that there is a small gap in the proof of \cite[Theorem~3]{TimurUniTri}. 
Indeed, Nasybullov's theorem uses \cite[Proposition~7]{TimurUniTri}, which in turn is cited as a version of Lev\v{c}uk's theorem for the case of integral domains and $n \geq 3$. However, as stated, \cite[Proposition~7]{TimurUniTri} does \emph{not} include automorphisms of type $\Sigma$ --- that is, the subgroup $\mc{U}^{(c)} \leq \Aut(\mbU_n(R))$ was not considered in~\cite{TimurUniTri}. The proposition thus cannot be applied as stated since there exist integral domains $R$ for which the subgroup $\mc{U}^{(c)} \leq \Aut(\mbU_n(R))$ is nontrivial --- take e.g., $R = \Z[1/2]$. Moreover when $n \leq 4$ the generators of $\Aut(\mbU_n(R))$ also include automorphisms not considered in~\cite{TimurUniTri}; cf. \cite[Theorems~2 and~3 and Corollary~5]{LevchukOriginal}.

Fortunately, the above mentioned gap in~\cite[Theorem~3]{TimurUniTri} can be overcome. First of all the assumptions on $R$ imply that it has characteristic zero and thus $n \geq 5$ since $R$ contains a copy of $\Z$. Secondly, the omitted automorphisms of type $\Sigma$ act trivially on almost all quotients of the lower central series of $\mbU_n(R)$. More precisely, any $\sigma \in \mc{U}^{(c)}$ induces the identity on $\mbU_n(R) / \gamma_{n-2}(\mbU_n(R))$, as can be seen from the description of $\mc{U}^{(c)}$; cf. Lemma~\ref{obs:sigmaautomorphisms}. The remaining arguments in the proof of \cite[Theorem~3]{TimurUniTri} thus carry over with the appropriate modifications. We take the opportunity to thank T.~Nasybullov for promptly discussing his results with us. 
\end{rmk}

\section{Applications of Theorem~\ref{thm:Additive}}\label{applications}

Here we exhibit new families of soluble matrix groups having property~$\Ri$ by applying Theorem~\ref{thm:Additive}. We keep the notation of the previous section for the automorphisms {$\addi{\alpha}$ and $\tau_\alpha$.}

\begin{pps} \label{pps:FptFptt-1OK} 
Let $R_0 \in \set{\Z, \F_p \mid p \in \N \text{ a prime}}$ and suppose $R \in \set{R_0[t], \laurent{R_0}, \ri_\K}$, where $\ri_\K$ is the ring of integers of an \emph{arbitrary} algebraic number field $\K$. Then $R(\addi{\alpha}) = \infty = R(\tau_\alpha)$ for any ring automorphism $\alpha \in \anel{\Aut}(R)$. In particular, the groups $\mbB_n(R)$ and $\PB_n(R)$ have property~$\Ri$ when $n \geq 4$.
\end{pps}

We split the proof in Sections~\ref{sec:Fpt}, \ref{sec:RinftyFptt-1} and~\ref{sec:OK} below. 
This result is an intermediate step towards Theorem~\ref{thm:Lieapplication}, covering the cases where $\RSpec(\Gamma_{n,p}) = \{\infty\}$ for $ n \geq 4$. 

\subsection{First case: polynomials on one variable \texorpdfstring{$R = R_0[t]$}{R = R0[t]}} \label{sec:Fpt} In this section we work with the polynomial rings $R = R_0[t]$ on one variable with coefficients in $R_0 \in \set{\Z, \F_p \mid p \text{ a prime}}$. 

It is a simple exercise to verify that every $\alpha \in \anel{\Aut}(R_0[t])$ is of the form $\alpha(\sum_{i=0}^d f_i t^i) = \sum_{i=0}^d f_i (at+b)^i$ for some $a \in R_0^\times$ and $b \in R_0$. We may assume $(a,b) \neq (1,0)$ for otherwise $\alpha$ would be the identity and obviously $R(\addi{\alpha})=\infty=R(\tau_\alpha)$ as $(R_0[t],+)$ is infinite abelian.

We start with the case $R_0 = \F_p$.

\begin{lem} \label{6.1exp1}
Let $h(t) \in \F_p[t]$. Given any $d_0 \geq \mathrm{deg}(h)$, the polynomial $h(t)$ can be uniquely written according to powers of $t$ modulo $p$ as follows. 
\[ h(t) = \sum_{\ell = 0}^{p-1} \sum_{\theta=0}^{d} h_{p\theta+\ell}t^{p\theta+\ell}, \]
where $d \in \Z_{\geq 0}$ is unique such that $d_0 = pd+r$ with $0 \leq r < p$.
\end{lem}

\begin{proof}
Writing $h(t) = \sum_{i=0}^{\mathrm{deg}(h)}h_i t^i$ uniquely and setting $h_{m} = 0$ whenever $m > \mathrm{deg}(h)$, simply rearrange the sum 
\begin{align*}
\sum_{i=0}^{\mathrm{deg}(h)}h_i t^i &= \sum_{i=0}^{d_0}h_i t^i = \sum_{i=0}^{pd+p-1}h_i t^i \, \text{ (here } h_i=0 \text{ when } i>\mathrm{deg}(h)\text{)} \\
&=\sum_{\theta=0}^{d}\left( h_{p\theta}t^{p\theta}+h_{p\theta+1}t^{p\theta+1}+\cdots+h_{p\theta+p-1}t^{p\theta+p-1} \right) \\ 
&= \sum_{\ell = 0}^{p-1} \sum_{\theta=0}^{d} h_{p\theta+\ell}t^{p\theta+\ell}.
\end{align*}
\end{proof}

\begin{lem} \label{claimFpt} Given any two polynomials $h(t), g(t) \in \F_p[t]$ and $d_0 \geq \mathrm{max}\{\mathrm{deg}(h),\mathrm{deg}(g)\}$, the polynomial $h(t) - \alpha(g(t))$ can be uniquely written as
\begin{align*}
h(t) - \alpha(g(t))=\sum_{\theta=1}^{d}( h_{p\theta+p-1} - g_{p\theta+p-1}a^{p\theta})t^{p\theta+p-1} + f_{h,g}(t),
\end{align*}
where $d \in \Z_{\geq 0}$ is unique such that $d_0 = pd+r$ with $0 \leq r < p$, and the remaining polynomial $f_{h,g}(t) \in \F_p[t]$ is such that \emph{no} power of $t$ appearing in $f_{h,g}(t)$ is of the form $t^{pk+p-1}$ with $k\geq 1$. 
\end{lem}

\begin{proof}
Notice that $\alpha$ preserves the degree of polynomials. Thus, applying Lemma~\ref{6.1exp1} to $h(t)$ and to $\alpha(g(t))$ yields
\[h(t)-\alpha(g(t)) = \sum_{\ell = 0}^{p-1} \sum_{\theta=0}^{d} h_{p\theta+\ell}t^{p\theta+\ell} - \sum_{\ell = 0}^{p-1} \sum_{\theta=0}^{d} g_{p\theta+\ell}(at+b)^{p\theta+\ell}. \]
(Recall that $h_i = 0$ (resp. $g_j = 0$) whenever $i > \mathrm{deg}(h)$ (resp. $j > \mathrm{deg}(g)$).) 
By isolating the sums with $\ell = p-1$ and $\theta \geq 1$, the summation above can be rewritten as
\begin{align} \label{supersum}
 \left(\sum_{\theta=1}^{d} (h_{p\theta+p-1}t^{p\theta+p-1}-g_{p\theta+p-1}(at+b)^{p\theta+p-1}) \right) + f_0(t),
\end{align}
where $f_0(t)$ is defined as
\begin{align*}
 f_0(t) =& \phantom{a} (h_{p-1}t^{p-1}-g_{p-1}(at+b)^{p-1}) \\
 &+ \sum_{\ell = 0}^{p-2} \sum_{\theta=0}^{d} (h_{p\theta+\ell}t^{p\theta+\ell} - g_{p\theta+\ell}(at+b)^{p\theta+\ell}).
\end{align*}

We now look at summands of the form $(at+b)^{p\theta+\ell}$ occurring in the polynomial $f_0(t)$. Since taking $p$-powers is an additive homomorphism in $\F_p[t]$, we have
\begin{align*} 
(at+b)^{p\theta+\ell} & = (at+b)^{p\theta}(at+b)^\ell \\
& = \begin{cases} a^{p\theta+\ell}t^{p\theta+\ell} & \text{if } b=0, \\ (a^{p\theta}t^{p\theta}+b^{p\theta}) \left( \sum_{m=0}^\ell \binom{\ell}{m} a^{\ell-m}b^m t^{\ell-m} \right) & \text{otherwise}. \end{cases}
\end{align*}
In particular, if $\ell < p-1$, then no power of $t$ appearing in the polynomial $(at+b)^{p\theta+\ell}$ is of the form $t^{pk+p-1}$ with $k \geq 1$. Hence the polynomial $f_0(t)$ defined above also has this property. 

Let us now consider the term $(at+b)^{p\theta+p-1}$ occurring in the summation~\bref{supersum}. Recalling that $x^{p-1} = 1$ for all $x \in \F_p^\times$, one has 
\begin{align*}
(at+b)^{p\theta+p-1} & = (at+b)^{p\theta}(at+b)^{p-1} \\
& = \begin{cases} a^{p\theta}t^{p\theta+p-1} \quad \text{ if } b=0, \\ 
\sum_{m=0}^{p-1}\binom{p-1}{m}\left( a^{p\theta-m} b^m t^{p\theta+p-1-m} + a^{-m}b^{p\theta+m}t^{p-1-m} \right). \end{cases} \\
& = a^{p\theta}t^{p\theta+p-1} + r_{\theta}(t),
\end{align*}
where 
\begin{align*}
r_{\theta}(t) = \begin{cases} 0 \quad \text{ if } b=0, \quad \text{ and otherwise} \\ 
b^{p\theta}t^{p-1}+\sum_{m=1}^{p-1}\binom{p-1}{m}\left( a^{p\theta-m} b^m t^{p\theta+p-1-m} + a^{-m}b^{p\theta+m}t^{p-1-m} \right). \end{cases}
\end{align*}
Therefore, in order for a power of $t$ of the form $t^{pk+p-1}$ with $k\geq 1$ to appear in the polynomial $(at+b)^{p\theta+p-1}$ above, one must have $b=0$ {and} $\theta \geq 1$, or in case $b\neq 0$ one must have $\theta \geq 1$ and $m=0$ in the second expression. In both cases, the only term where the given power occurs is precisely $a^{p\theta}t^{p\theta+p-1}$. Put differently, for any $\theta$ the polynomial $r_\theta(t)$ is such that \emph{no} power of $t$ occurring in it is of the form $t^{pk+p-1}$ with $k\geq 1$ because $1 \leq m \leq p-1$.

Now define a new polynomial $f_{h,g}(t) \in \F_p[t]$ by  
\[f_{h,g}(t) := f_0(t) - \sum_{\theta=1}^d r_{\theta}(t).\]
Since neither $f_0(t)$ nor the $r_{\theta}(t)$ contain $t$-powers of the form $t^{pk+p-1}$ with $k \geq 1$, the same holds for $f_{h,g}(t)$ by definition. Expanding the summation~\bref{supersum}, one readily checks that 
\begin{align*}
 h(t) - \alpha(g(t)) & = \sum_{\theta=1}^d (h_{p\theta+p-1}-g_{p\theta+p-1}a^{p\theta})t^{p\theta+p-1} - \sum_{\theta=1}^d r_{\theta}(t) + f_0(t) \\
 & = \sum_{\theta=1}^d (h_{p\theta+p-1}-g_{p\theta+p-1}a^{p\theta})t^{p\theta+p-1} + f_{h,g}(t),
\end{align*}
as desired.  
\end{proof}

We now proceed with the proof of Proposition~\ref{pps:FptFptt-1OK} in the present case. Consider the family of polynomials $\set{t^{p(p-1)i+p-1}}_{i \in \N} \subset R_0[t] = \F_p[t]$. Given natural numbers $i > j \geq 1$, one has that $t^{p(p-1)i+p-1}$ is $\addi{\alpha}$-conjugate to $t^{p(p-1)j+p-1}$ if and only if there exists $h(t) \in R_0[t]$ of degree $\mathrm{deg}(h) \geq {p(p-1)i+p-1}$ such that 
\[ t^{p(p-1)i+p-1} - t^{p(p-1)j+p-1} = h(t) - \alpha(h(t)). \]
Applying Lemma~\ref{claimFpt} with $g(t)=h(t)$ and $d_0 \geq \mathrm{deg}(h)$, the above holds if and only if 
\[ t^{p(p-1)i+p-1} - t^{p(p-1)j+p-1} = \left(\sum_{\theta=1}^{d} h_{p\theta+p-1}(1-a^{p\theta})t^{p\theta+p-1}\right) + f_{h,h}(t), \]
where $f_{h,h}(t)$ and $d$ are as in the statement of Lemma~\ref{claimFpt}. Since $f_{h,h}(t)$ contains no $t$-powers of the form $pk+p-1$ with $k\geq 1$, it follows that some $\theta \in \set{1, \ldots, d}$ must equal $(p-1)i$, yielding an equality for coefficients of the term $t^{p(p-1)i+p-1}$. But one would then obtain
\[ 1 = h_{p(p-1)i+p-1}(1-a^{p(p-1)i}) = h_{p(p-1)i+p-1}(1-(a^{p-1})^{pi}) = 0 \]
because $a^{p-1}=1$. In other words such a polynomial $h(t)$ cannot exist, i.e., $t^{p(p-1)i+p-1}$ and  $t^{p(p-1)j+p-1}$ define disjoint $\addi{\alpha}$-twist conjugacy classes whenever $i > j$. Thus $R(\addi{\alpha}) = \infty$.

Showing that $R(\tau_\alpha) = \infty$ is similar to the previous case. Indeed, we shall verify that the elements $(t^{p(p-1)i+p-1},0)$ and $(0,-t^{p(p-1)j+p-1}) \in \F_p[t] \times \F_p[t]$, $i,j\in \N$, belong to distinct $\tau_\alpha$-twist conjugacy classes whenever $i > j$. 

By definition, $(t^{p(p-1)i+p-1},0)$ and $(0,-t^{p(p-1)j+p-1})$ are $\tau_\alpha$-conjugate if there exist $h(t), g(t) \in \F_p[t]$ of degrees $\mathrm{deg}(h),\mathrm{deg}(g) \geq {p(p-1)i+p-1}$ such that 
\begin{align*} 
(t^{p(p-1)i+p-1}, \, \, t^{p(p-1)j+p-1}) & = (h(t), \, \, g(t)) - \tau_\alpha( \, (h(t), \, g(t)) \, ) \\
 & = (\, h(t) - \addi{\alpha}(g(t)), \, \, g(t) - \addi{\alpha}(h(t)) \, ).
\end{align*}
Again applying Lemma~\ref{claimFpt} twice with $d_0 \geq \mathrm{max}\{\mathrm{deg}(h),\mathrm{deg}(g)\}$, the above holds if and only if 
\begin{align*} 
& (t^{p(p-1)i+p-1}, \quad t^{p(p-1)j+p-1}) \\
=& \left( \sum_{\theta=1}^{d} (h_{p\theta+p-1} - g_{p\theta+p-1} a^{p\theta}) t^{p\theta+p-1} + f_{h,g}(t) , \right. \\ 
& \, \, \left. \sum_{\theta=1}^{d} (g_{p\theta+p-1} -  h_{p\theta+p-1} a^{p\theta}) t^{p\theta+p-1} + f_{g,h}(t) \right).
\end{align*}
Since the polynomials $f_{h,g}(t)$ and $f_{g,h}(t)$ from Lemma~\ref{claimFpt} contain no $t$-powers whose exponents are of the form $pk+p-1$ with $k \geq 1$, we have that some $\theta \in \set{1,\ldots,d}$ must equal $(p-1)i$. This yields the system of coefficients
\[ \begin{cases}
 1 = h_{p(p-1)i+p-1} - g_{p(p-1)i+p-1}a^{p(p-1)i} = h_{p(p-1)i+p-1} - g_{p(p-1)i+p-1}, \\
 0 = g_{p(p-1)i+p-1} - h_{p(p-1)i+p-1}a^{p(p-1)i} = g_{p(p-1)i+p-1} - h_{p(p-1)i+p-1}, 
\end{cases}\]
which has no solutions over $R_0 = \F_p$. Thus the existence of polynomials $h(t), g(t) \in R_0[t]$ satisfying the above is impossible, whence $(t^{p(p-1)i+p-1},0)$ and $(0, -t^{p(p-1)j+p-1})$ lie in distinct $\tau_\alpha$-twist conjugacy classes whenever $i > j$. The desired result on property $\Ri$ now follows from Theorem~\ref{thm:Additive}. 

The case $R_0 = \Z$ follows from the previous one. Take, for instance, the canonical projection $\pi : \Z \onto \Z / 2\Z = \F_2$. Given $\alpha \in \anel{\Aut}(\Z[t])$, which acts via $\alpha(\sum_{i=0}^d f_i t^i) = \sum_{i=0}^d f_i (at+b)^i$, it is clear that the ideal $2\Z[t]$ --- i.e. polynomials with even coefficients --- is $\addi{\alpha}$-invariant. Thus $\addi{\alpha}$ induces the automorphism $\barra{\addi{\alpha}}$ on $\Z[t] / 2\Z[t] = \F_2[t]$ given by $\addi{\alpha}(\sum_{i=0}^d \pi(f_i) t^i) = \sum_{i=0}^d \pi(f_i) (\pi(a)t + \pi(b))^i$. But $R(\addi{\alpha}) \geq R(\barra{\addi{\alpha}})$ by Lemma~\ref{lem:desempre}, and we have just shown that such an $\barra{\addi{\alpha}}$ on $\F_2[t]$ has infinite Reidemeister number. Similarly, the product of ideals $2\Z[t] \times 2\Z[t]$ is also $\tau_\alpha$-invariant, whence taking the induced automorphism $\barra{\tau}_\alpha$ on $\F_2[t] \times \F_2[t]$ yields $R(\tau_\alpha) \geq R(\barra{\tau}_\alpha) = \infty$ by the previous considerations. This finishes off the proof of the present case. \qed

\subsection{The case of Laurent polynomials \texorpdfstring{$R = \laurent{R_0}$}{R = R0[t,1/t]}} \label{sec:RinftyFptt-1} Turning to the second claim of Proposition~\ref{pps:FptFptt-1OK} we let $R$ denote throughout this section the ring of Laurent polynomials $R = R_0[t,t^{-1}]$. Here, $R_0$ is either the ring of integers $\Z$ or a finite field $\F_p$ with a prime number $p$ of elements.

Straightforward computations show that the group of ring automorphisms $\anel{\Aut}(\laurent{R_0})$, regardless of $R_0 \in \{\F_p,\Z\}$, is just the cyclic group of order two, generated by the automorphism which takes $t$ to $t^{-1}$. Thus, given $\alpha \in \anel{\Aut}(\laurent{R_0})$ and an arbitrary Laurent polynomial $f(t) = \sum_{\ell=-\infty}^\infty f_\ell t^\ell \in \laurent{R_0}$, where only finitely many coefficients $f_\ell\in R_0$ are non-zero, we have $\alpha(f(t)) = \sum_{\ell=-\infty}^\infty f_\ell t^{\pm \ell}$. 

Similarly to the case of $R_0[t]$ from Section~\ref{sec:Fpt}, we prove that the elements $t^\ell \in \laurent{R_0}$ yield infinitely many $\addi{\alpha}$-twist conjugacy classes. Take $i > j \in \N$. Then $t^i$ and $t^j$ are $\addi{\alpha}$-conjugate if and only if there exists $h(t) = \sum_{\ell=-\infty}^\infty h_\ell t^\ell \in \laurent{R_0}$ such that $0\neq t^i - t^j = (\id-\addi{\alpha})(h(t))$. If $\alpha$ sends $t$ to $t$, then $\addi{\alpha}$ is itself the identity, whence there is no such $h(t)$ and $R(\addi{\alpha}) = \infty$. Assume otherwise that $\addi{\alpha}(h(t)) =  \sum_{\ell=-\infty}^\infty h_\ell t^{-\ell}$. Since $i$ is positive and distinct from $j$, we then have the system 
\[ \begin{cases}
 h_i - h_{-i} = 1, \\
 -h_i + h_{-i} = 0
\end{cases}\]
of coefficients, which has no solutions over $R_0$. Thus such an element $h(t) \in \laurent{R_0}$ never exists, that is $t^i$ and $t^j$ can not be $\addi{\alpha}$-conjugate if $i \neq j$, whence $R(\addi{\alpha}) = \infty$.

We now check that the additive automorphism $\tau_\alpha((r,s)) = (\addi{\alpha}(s),\addi{\alpha}(r))$ of $\laurent{R_0} \times \laurent{R_0}$ also has $R(\tau_\alpha) = \infty$. Similarly to the above we show that the elements $(t^i,0)$, $(0,-t^j) \in \laurent{R_0} \times \laurent{R_0}$ with $i > j \in \N$ define distinct $\tau_\alpha$-twist conjugacy classes, which in turn implies $R(\tau_\alpha) = \infty$. Again, $(t^i,0)$ and $(0,-t^j)$ are $\tau_\alpha$-conjugate if and only if there exists a pair of Laurent polynomials $(h(t), g(t)) = (\sum_{\ell=-\infty}^\infty h_\ell t^\ell, \sum_{m=-\infty}^\infty g_m t^m) \in R_0[t,t^{-1}] \times R_0[t,t^{-1}]$ satisfying 
\begin{align*} 
(t^i, \, \, t^j) & = (h(t), \, \, g(t)) - \tau_\alpha( \, (h(t), \, g(t)) \, ) \\
 & = (\, h(t) - \addi{\alpha}(g(t)), \, \, g(t) - \addi{\alpha}(h(t)) \, )
\end{align*}
\[ = \left( \sum_{\ell=-\infty}^\infty h_\ell t^\ell - \sum_{m=-\infty}^\infty g_m t^{\pm m} , \, \, \sum_{m=-\infty}^\infty g_m t^m - \sum_{\ell=-\infty}^\infty h_\ell t^{\pm \ell} \right).
\]
One has the following two possible systems of coefficients. Firstly, if $\alpha$ sends $t$ to $t$, we have the system 
\[ \begin{cases}
 1 = h_i - g_i, \\
 0 = g_i - h_i. 
\end{cases}\]
Secondly, if $\alpha$ sends $t$ to $t^{-1}$, the system we obtain is 
\[ \begin{cases}
 1 = h_i - g_{-i}, \\
 0 = g_{-i} - h_i. 
\end{cases}\]
In both cases, no solutions exist over $R_0$, so that there is no such pair $(h(t), g(t))$ satisfying the above. Therefore $(t^i,0)$ and $(0,-t^j)$ lie in distinct $\tau_\alpha$-twist conjugacy classes for all $i > j \in \N$, yielding $R(\tau_\alpha) = \infty$. That the groups $\mbB_n(\laurent{R_0})$ and $\PB_n(\laurent{R_0})$ have $\Ri$ for $n\geq 4$ now follows immediately from Theorem~\ref{thm:Additive}. \qed

\subsection{The case where \texorpdfstring{$R = \ri_\K$}{R = Ok}, a ring of integers} \label{sec:OK} We now prove the last claim of Proposition~\ref{pps:FptFptt-1OK}. 
Recall that the integral closure of $\Z$ in a finite extension of $\Q$ is a free $\Z$-module of finite rank, whence the underlying additive group $\Addi(\ri_\K)$ of any ring of integers $\ri_\K$ is finitely generated. 

Suppose $\alpha$ is any ring automorphism of $\ri_\K$. To prove the claim, it suffices to check that $\addi{\alpha}$ and $\tau_\alpha$ have infinitely many fixed points due to Lemma~\ref{lem:ReidemeisterAbelian} and the observation above. This is obviously true for $\addi{\alpha}$ since $\Z \subseteq \ri_\K$ is always contained in the set of fixed points of $\alpha$. And for the automorphism $\tau_\alpha : \ri_\K \times \ri_\K \to \ri_\K \times \ri_\K$, $\tau_\alpha(r,s) = (\addi{\alpha}(s), \, \addi{\alpha}(r))$, the argument is the same. Indeed, the diagonal of $\Z \times \Z$ is clearly contained in the set of fixed points of $\tau_\alpha$. Using Theorem~\ref{thm:Additive} we conclude the proof of Proposition~\ref{pps:FptFptt-1OK}. \qed

\section{Some groups in positive characteristic that do not have property \texorpdfstring{$\Ri$}{R infinity}} \label{sec:newexampleswithoutRinfty}

In contrast with Theorem~\ref{thm:Additive} we now give examples of metabelian $S$-arithmetic groups in positive characteristic not having property~$\Ri$. Our list includes finitely generated and non-finitely generated examples.

\begin{pps}\label{nori} Let $\F_q$ denote the finite field with $q$ elements. 
\begin{enumerate}
	\item \label{nori.i} 
	For every irreducible monic polynomial $P(X) \in \F_q[X]$ of degree at least $2$, there exist automorphisms 
	\[\phee_{\mathbb{A},P} \in \Aut(\PB_2(\F_q[t])), \,\, \phee_{\mbB,P} \in \Aut(\mbB_2(\F_q[t])), \,\, \Phi_P \in \Aut(\mbU_2(\F_q[t]))\]
	with Reidemeister numbers 
    \[ R(\phee_{\mathbb{A},P}) = q-1, \quad R(\phee_{\mbB,P}) = (q-1)^2, \quad R(\Phi_P) = 1, \]
    respectively. In particular, 
	\emph{none} of the groups above have property~$\Ri$.  
	\item \label{nori.ii} 
	If $q \geq 4$, there exist automorphisms $\phee_{A} \in \Aut(\PBplus_2(\F_q[t,t^{-1}]))$, $\phee_{B} \in \Aut(B^+_2(\F_q[t,t^{-1}]))$ and $\phee' \in \Aut(\mbU_2(\F_q[t,t^{-1}]))$ 
	with 
    \[R(\phee_A) = 2, \quad R(\phee_B) = 4, \quad R(\phee') = 1,\] respectively. In particular, 
	\emph{none} of the groups above have property~$\Ri$.
\end{enumerate}
\end{pps} 

The groups $\PBplus_n(R)$ and $B_n^+(R)$ are just variants of $\PB_n(R)$ and $\mbB_n(R)$, respectively, without torsion on the diagonal part; cf. Section~\ref{fqttmo} for the definition. 

Proposition~\ref{nori} has some overlap with the existing literature. For $q = p$ a prime, the group $\PBplus_2(\F_p[t,t^{-1}])$ is isomorphic to the lamplighter group $\mc{L}_p = C_p \wr \Z$. 
Using different techniques, Gon\c{c}alves and Wong~\cite{DacibergWongWreath} completely classified which wreath products of the form $A \wr \Z$ have property~$\Ri$, where $A$ is a finitely generated abelian group. In particular, for the small fields $\F_2$, $\F_3$ excluded from Proposition~\dref{nori}{nori.ii}, their results imply that $\PBplus_2(\laurent{\F_2})$ and $\PBplus_2(\laurent{\F_3})$ actually \emph{do have} property~$\Ri$, while $\PBplus_2(\laurent{\F_q})$ does not when $q\geq 4$.

Combining Propositions~\ref{pps:FptFptt-1OK} and~\ref{nori} yields the following. 
\[
\begin{array}{c|c|c|c} 
\text{Does } G_n(R) \text{ have } \Ri ? & n = 2 & n = 3 & n \geq 4 \\ 
\hline 
R = \F_q[t] & \text{No} & \text{Unknown}	&	\text{Yes} \\ 
G_n = \mbB_n, \PB_n &  & & \\
\hline  
R = \F_q[t,t^{-1}] & \text{No, if } q \geq 4, & \text{Unknown}	&	\text{Yes, if} \\ 
& G_n = B_n^+, \PBplus_n & & G_n = \mbB_n, \PB_n \\
\hline
\end{array}
\]
It is unclear whether the groups in the $3\times 3$ case should have property~$\Ri$. The `plus' versions $B_n^+(\F_q[t])$ and $\PBplus_n(\F_q[t])$ are isomorphic to $\mbU_n(\F_q[t])$ by definition, and these groups do not have~$\Ri$ whenever $n < q$; see \cite[Proposition~8]{TimurUniTri}. Whether $B_n^+(\laurent{\F_q})$ and $\PBplus_n(\laurent{\F_q})$ have~$\Ri$ for large~$n$, and whether $\mbB_2(\laurent{\F_q})$ and $\PB_2(\laurent{\F_q})$ do not have~$\Ri$, is also unknown to us. 

\begin{rmk}
We stress that there cannot be a version of Theorem~\ref{thm:Additive} in the `plus' case removing torsion from the diagonal---at least not for all $n \geq 4$. That is, the hypotheses of Theorem~\ref{thm:Additive} might hold while $B_n^+(R)$ and $\PBplus_n(R)$ do not necessarily have~$\Ri$. For example, take $R = \ri_\K$ a ring of integers of a number field $\K$ such that $\ri_\K$ has finitely many units. As in Section~\ref{sec:OK}, the maps $\alpha_{\mathrm{add}}$ and $\tau_\alpha$ always have infinite Reidemeister number over $\ri_\K$. On the other hand, $B_n^+(\ri_\K) \cong \PBplus_n(\ri_\K) \cong \mbU_n(\ri_\K)$ and, by \cite[Proposition~8]{TimurUniTri}, it is known that $\RSpec(\mbU_4(\ri_\K)) \neq \{\infty\}$. 
\end{rmk}

\subsection{Proof of Proposition~\dref{nori}{nori.i}}\label{fqt}
Throughout this section, $R$ denotes $\F_q[t]$, the polynomial ring on one variable over $\F_q$. 

The fact that there exist monic polynomials $P(X) = a_0 + a_1 X +  \ldots + a_{d-1} X^{d-1} + X^d $ of degree $d \geq 2$ with coefficients in $\F_q$ which are \emph{irreducible} over $\F_q$ is straightforward from the fact that $\F_q$ is not algebraically closed. Now fix such a polynomial $P(X)$. Certainly $a_0 \neq 0$ for otherwise $P(X)$ would be divisible by $X$. The companion matrix 
\[ C_P = \begin{pmatrix} 0 & 0 & \ldots & 0 & -a_0 \\ 1 & 0 & \ldots & 0 & -a_1 \\ 0 & 1 & \ldots & 0 & -a_2 \\ \vdots & \vdots & \ddots & \vdots & \vdots \\ 0 & 0 & \ldots & 1 & -a_{d-1} \end{pmatrix} \]
of $P(X)$ thus has determinant $\det(C_P) = a_0 \neq 0$ and therefore defines an element, also denoted $C_P$ by abuse of notation, of $\GL_{d}(\F_q)$. Moreover, the characteristic polynomial of $C_P$ equals $P(X)$ itself. Since $P(X)$ is irreducible, one has {$\det(C_P - \lambda\cdot\mb{1}_d) \neq 0$} for all $\lambda \in \F_q \backslash \set{0}$, hence
\begin{align} \label{eq:Lula}
{\det(\id - a \cdot C_P) \neq 0} 
%\det(\mb{1}_d-a \cdot C_P) \neq 0 
\text{ for any scalar } a \in \F_q \backslash \set{0}.
\end{align}
That is, the linear map $\id - a \cdot C_P$ lies in $\GL_d(\F_q)$ if $a \neq 0$.

We would like to define an additive automorphism $\Phi_P \in \Aut(\Addi(\F_q[t]))$ induced by $C_P$. To do this we first write $\F_q[t]$ as an infinite dimensional $\F_q$-vector space on the standard basis $\set{1, t, t^2, \ldots}$, that is 
\[ \F_q[t] = \bigoplus_{\ell=0}^\infty \F_q \cdot t^\ell. \]
We remark that the elements $x \in \F_q \subset \F_q[t]$ can be viewed both as vectors $xt^0+0t^1+0t^2+\ldots \in \oplus_{\ell=0}^\infty \F_q t^\ell$ and as scalars of the ground field $\F_q$. In particular, an $\F_q$-linear transformation $T : \F_q[t] \to \F_q[t]$ \emph{needs not} fix the {vector} $xt^0+0t^1+0t^2+\ldots$ corresponding to the element $x \in \F_q$, though it still satisfies $T(x \cdot v) = x \cdot T(v)$ for any vector $v \in \F_q[t]$ and scalar $x \in \F_q$.

To construct our automorphism $\Phi_P$ we decompose $\F_q[t]$ into finite-dimensional blocks. More precisely, an arbitrary vector $v \in \F_q[t]$ shall be written as a sum $v = v_1+v_2+v_3+\ldots$ where each $v_k$ lies in the $(k+1)$-st $d$-dimensional block of the space $\oplus_{\ell=0}^\infty \F_q t^\ell$. That is, we decompose 
\[ \F_q[t] = \bigoplus_{\ell=0}^\infty \F_q \cdot t^\ell = \bigoplus_{k=0}^\infty \bigoplus_{\ell=kd}^{(k+1)d-1} \F_q \cdot t^\ell \]
and write $v \in \F_q[t]$ as
\[ v = \sum_{k=0}^\infty v_k \, \text{ where } \, v_k \in \bigoplus_{\ell=kd}^{(k+1)d-1} \F_q \cdot t^\ell \cong \F_q^d. \]
Since $C_P \in \GL_d(\F_q)$, we may define on each $d$-dimensional block $\oplus_{\ell=kd}^{(k+1)d-1} \F_q t^\ell \cong \F_q^d$ above the (invertible) linear transformation 
\[ \xymatrix@R=2mm{ 
\displaystyle C_{P,k} : \bigoplus_{\ell=kd}^{(k+1)d-1} \F_q \cdot t^\ell \ar[r] & \displaystyle\bigoplus_{\ell=kd}^{(k+1)d-1} \F_q \cdot t^\ell \\
x \ar@{|->}[r] & C_P(x)
}
 \]
which is just a copy of $C_P$. Thus the map $\Phi_P : \F_q[t] \to \F_q[t]$ defined via 
\[ 
\xymatrix@R=2mm{ 
\Phi_P : \displaystyle\bigoplus_{k=0}^\infty \bigoplus_{\ell=kd}^{(k+1)d-1} \F_q \cdot t^\ell \ar[r] & \displaystyle \bigoplus_{k=0}^\infty \bigoplus_{\ell=kd}^{(k+1)d-1} \F_q \cdot t^\ell \\
v = \displaystyle\sum_{k=0}^\infty v_k \ar@{|->}[r] & \displaystyle\sum_{k=0}^\infty C_{P,k}(v_k)
} 
\]
is an automorphism of $\F_q[t]$ as an infinite-dimensional $\F_q$-vector space because each $C_{P,k}$ is an invertible linear transformation of the $d$-dimensional subspace $\oplus_{\ell=kd}^{(k+1)d-1} \F_q t^\ell$. In particular, $\Phi_P \in \Aut(\Addi(\F_q[t]))$.

Given a unit $a \in \F_q \backslash \set{0}$, let $\mc{m}_a : \Addi(\F_q[t]) \to \Addi(\F_q[t])$ denote the additive automorphism which is just multiplication by $a$. We claim that 
\begin{equation} \label{eq:R(maPhi)}
R(\mc{m}_a \circ \Phi_P) = 1 \text{ for all } a \in \F_q \backslash \set{0},
\end{equation} 
where $\mc{m}_a$ and $\Phi_P$ are here viewed as automorphisms of $\Addi(\F_q[t])$. Let then $r = \sum_{k=0}^\infty r_k \in \F_q[t]$ be arbitrary, where the $r_k$ are as above. Since $\id-a\cdot C_{P,k} \in \GL_d(\F_q)$ by the definition of $C_{P,k}$ and by equality~\bref{eq:Lula}, we may define 
\[ s_k := (\id-a\cdot C_{P,k})^{-1}(r_k) \in \bigoplus_{\ell=kd}^{(k+1)d-1} \F_q \cdot t^\ell \]
for all $k$ and set $s = \sum_{k=0}^\infty s_k \in \F_q[t]$. One then has
\begin{align*} 
r = \sum_{k=0}^\infty r_k = \sum_{k=0}^\infty (\id-a\cdot C_{P,k})(s_k) & = (\id-a\cdot\Phi_P)\left(\sum_{k=0}^\infty s_k\right) \\ & = s - \mc{m}_a \circ \Phi_P(s).
\end{align*}
Thus $r$ is $(\mc{m}_a\circ\Phi_P)$-conjugate to the zero vector $0 \in \F_q[t]$, whence the claim.

Equality~\bref{eq:R(maPhi)} alone implies that $\mbU_2(\F_q[t]) = \left( \begin{smallmatrix} 1 & \F_q[t] \\ 0 & 1 \end{smallmatrix} \right) \cong \Addi(\F_q[t])$ does not have property~$\Ri$. To verify that the same is true for the groups $\PB_2(\F_q[t])$ and $\mbB_2(\F_q[t])$, consider the maps 
\[
\xymatrix@R=2mm{ \phee_{\mathbb{A},P} : \PB_2(\F_q[t]) \ar[r] & \PB_2(\F_q[t]) \\ 
{\begin{bmatrix} u & r \\ 0 & v \end{bmatrix}} \ar@{|->}[r] & {\begin{bmatrix} u & \Phi_P(r) \\ 0 & v \end{bmatrix}} } 
\]
and 
\[ 
\xymatrix@R=2mm{ \phee_{\mathbf{B},P} : \mbB_2(\F_q[t]) \ar[r] & \mbB_2(\F_q[t]) \\ 
{\begin{pmatrix} u & r \\ 0 & v \end{pmatrix}} \ar@{|->}[r] & {\begin{pmatrix} u & \Phi_P(r) \\ 0 & v \end{pmatrix}}.} 
\]
These are in fact automorphisms because $\Phi_P : \F_q[t] \to \F_q[t]$ is $\F_q$-linear. By definition, $\phee_{\mathbb{A},P}$ and $\phee_{\mathbf{B},P}$ preserve the diagonal and unipotent parts of the given groups. 
Recalling that $\mbU_2(\F_q[t])$ is characteristic in the groups above and because their diagonal parts are finite, 
we know from Lemma~\ref{lem:desempre} that $R(\phee_{\mathbb{A},P})$ and $R(\phee_{\mathbf{B},P})$ are finite in case $R(\iota_{\mathbb{A}} \circ \phee_{\mathbb{A},P}')$ and $R(\iota_{\mathbf{B}} \circ \phee_{\mathbf{B},P}')$ are so for all $\iota_{\mathbb{A}} \in \Inn(\PB_2(\F_q[t]))$ and $\iota_{\mbB} \in \Inn(\mbB_2(\F_q[t]))$. Here, $\phee_{\mathbb{A},P}'$ and $\phee_{\mbB,P}'$ denote the restrictions to $\mbU_2(\F_q[t]) \cong \Addi(\F_q[t])$ of $\phee_{\mathbb{A},P}$ and $\phee_{\mbB,P}$, respectively. But direct matrix computations show that $\iota_{\mathbb{A}} \circ \phee_{\mathbb{A},P}'$ and $\iota_{\mathbf{B}} \circ \phee_{\mathbf{B},P}'$ are of the form $\mc{m}_a \circ \Phi_P$ and $\mc{m}_b \circ \Phi_P$, respectively, for some units $a,b \in \F_q \backslash \set{0}$. Since $R(\mc{m}_a \circ \Phi_P) = R(\mc{m}_b \circ \Phi_P) = 1$ by the equality~\bref{eq:R(maPhi)}, it follows that $\phee_{\mathbb{A},P}$ and $\phee_{\mbB,P}$ have finite Reidemeister number as well. Therefore $\PB_2(\F_q[t])$ and $\mbB_2(\F_q[t])$ do not have property $\Ri$.

Lastly, the Reidemeister numbers of $\phee_{\mathbb{A},P}$ and $\phee_{\mbB,P}$ can be promptly computed. Since both such automorphisms fix the diagonal pointwise, every diagonal matrix lies in a distinct Reidemeister class, so that $R(\phee_{\mathbb{A},P}) \geq q-1$ and $R(\phee_{\mbB,P}) \geq (q-1)^2$. Now, because every $r \in \F_q[t]$ can be written as $r = s - u \cdot \Phi_P(s)$ for any unit $u \in \F_q \backslash \{0\}$ and for some $s \in \F_q[t]$, it follows that an arbitrary element $\left( \begin{smallmatrix} u & r \\ 0 & v\end{smallmatrix} \right) \in \mbB_2(\F_q[t])$ is $\phee_{\mbB,P}$-conjugate to its diagonal part $\left( \begin{smallmatrix} u & 0 \\ 0 & v\end{smallmatrix} \right)$ since
\[ \begin{pmatrix} u & r \\ 0 & v \end{pmatrix} = \begin{pmatrix} 1 & s \\ 0 & v \end{pmatrix} \cdot \begin{pmatrix} u & 0 \\ 0 & v \end{pmatrix} \cdot \begin{pmatrix} 1 & \Phi_P(s) \\ 0 & v \end{pmatrix}^{-1}.\] 
(Analogously for $ \PB_2(\F_q[t])$.) Thus $R(\phee_{\mathbb{A},P}) = q-1$ and $R(\phee_{\mbB,P}) = (q-1)^2$, as claimed. \qed


\subsection{Proof of Proposition~\dref{nori}{nori.ii}} \label{fqttmo}
Throughout this section, $R$ denotes the Laurent polynomial ring $\F_q[t,t^{-1}]$ with $q\geq 4$ a power of a prime. 

The group of units of $R$ is 
\[R^{\times}=\{u t^k \mid u \in \F_q^\times, \, k \in \Z\},\] 
which has torsion subgroup $t(R^\times)=\F_{q}^{\times}$. The group $R^\times_{\mathrm{tf}}=\{t^k\mid k \in \Z\}$ is a torsion-free complement of $t(R^\times)$, and we define 
\begin{alignat*}{3} B_{n}^{+}(R)&:= \mbB_n(R ; R^\times_{\mathrm{tf}}) &&:= \left\{ b=(b_{i,j}) \in \mbB_n(R) \bigm| b_{i,i} \in R^\times_{\mathrm{tf}},\, b_{i,j} \in R \right\},\\	
\PBplus_n(R) &:= \PB_n(R ; R^\times_{\mathrm{tf}}) &&:= \left\{ b=[b_{i,j}] \in \PB_n(R) \bigm| b_{i,i} \in R^\times_{\mathrm{tf}},\, b_{i,j} \in R \right\}.
\end{alignat*}
In other words, $B_n^+(R)\leq \mbB_n(R)$ and $\PBplus_n(R) \leq \PB_n(R)$ are the subgroups whose nontrivial diagonal entries are torsion-free units. (These subgroups are not unique setwise since they depend on the choice of $R^\times_{\mathrm{tf}}$, but they are unique up to isomorphism.) 


\begin{lem}\label{lemmaa} Let $\F_q$ denote a finite field with $q \geq 4$ elements. There exists a nonzero $a \in \F_q$ such that $1-a^2$ is a unit of $\F_q$. In particular, $1-a$ is a unit of $\F_q$, and the system
\begin{align*} X-aY&=\ell\\
			  -aX+Y&=m
\end{align*}
has a solution $(X,Y) \in \F_{q}^{2}$ for all choices of $\ell, m \in \F_q$.
\end{lem}
\begin{proof} Take $a$ to be a generator of $\F_{q}^{\times}$. Since $|\F_q|\geq 4$, it follows that $a^2\neq 1$.  
\end{proof}

Let $a\in \F_{q}^{\times}$ be as in Lemma~\ref{lemmaa}. Let also $k,j \in \Z$ and $f(t) \in \F_q[t,t^{-1}]$. Define 
\[ 
\xymatrix@R=2mm{ \phee_{B} : B_2^+(\F_q[t,t^{-1}]) \ar[r] & B_2^+(\F_q[t,t^{-1}]) \\ 
{\begin{pmatrix} t^k & f(t) \\ 0 & t^j \end{pmatrix}} \ar@{|->}[r] & {\begin{pmatrix} t^{-k} & af(t^{-1}) \\ 0 & t^{-j} \end{pmatrix}}} 
\]
and 
\[ 
\xymatrix@R=2mm{ \phee_{A} : \PBplus_2(\F_q[t,t^{-1}]) \ar[r] & \PBplus_2(\F_q[t,t^{-1}]) \\ 
{\begin{bmatrix} t^k & f(t) \\ 0 & t^j \end{bmatrix}} \ar@{|->}[r] & {\begin{bmatrix} t^{-k} & af(t^{-1}) \\ 0 & t^{-j} \end{bmatrix}}.} 
\]
It is a simple exercise to check that these maps are automorphisms. As $\mbU_2(R)$ is characteristic in $B_2^+(R)$ and $\PBplus_2(R)$, we also consider the automorphism  $\phee'= \phee_{B}|_{\mbU_2(R)}=\phee_{A}|_{\mbU_2(R)}$ on $\mbU_2(R)$.

\begin{pps}\label{pps:norifqttmo} The automorphisms defined above satisfy 
\begin{align*}\mathcal{R}(\phee_{B})&=\left\{\left[\left(\begin{smallmatrix} 1&0\\0&1 \end{smallmatrix}\right)\right]_{\phee_{B}},\left[\left(\begin{smallmatrix} t&0\\0&1 \end{smallmatrix}\right)\right]_{\phee_{B}},\left[\left(\begin{smallmatrix} 1&0\\0&t \end{smallmatrix}\right)\right]_{\phee_{B}},\left[\left(\begin{smallmatrix} t&0\\0&t \end{smallmatrix}\right)\right]_{\phee_{B}}\right\},\\
\mathcal{R}(\phee_{A})&=\left\{\left[\left(\begin{smallmatrix} 1&0\\0&1 \end{smallmatrix}\right)\right]_{\phee_{A}},\left[\left(\begin{smallmatrix} t&0\\0&1 \end{smallmatrix}\right)\right]_{\phee_{A}}\right\} \text{ and }\\
\mathcal{R}(\phee')&=\left\{\left[\left(\begin{smallmatrix} 1&0\\0&1 \end{smallmatrix}\right)\right]_{\phee'}\right\},
\end{align*}
where $\mathcal{R}(\phi)$ denotes the set of Reidemeister classes of the automorphism $\phi$.
\end{pps}

\begin{proof} We show the result for $B_{2}^{+}(R)$, the proofs for $\PBplus_2(R)$ and for $\mbU_2(R)$ follow from entirely analogous arguments. 

Given elements of $B_{2}^{+}(R)$ of the form 
\[\left(\begin{smallmatrix} t^{k}&f(t)\\0&t^{\ell}\end{smallmatrix}\right)\text{ and }\left(\begin{smallmatrix} t^{x}&0\phantom{|}\\0&t^{y}\end{smallmatrix}\right),\] 
we have that 
\[ \left(\begin{smallmatrix} t^{k}&f(t)\\0&t^{\ell}\end{smallmatrix}\right)\left(\begin{smallmatrix} t^{x}&0\phantom{|}\\0&t^{y}\end{smallmatrix}\right)\phee_{B}\left(\left(\begin{smallmatrix} t^{k}&f(t)\\0&t^{\ell}\end{smallmatrix}\right)\right)^{-1} = \left(\begin{smallmatrix} t^{2k+x}&t^{\ell+y}f(t)-at^{2k+\ell+x}f(t^{-1})\\0&t^{2\ell+y}\end{smallmatrix}\right).\]

Consequently, an element $b=\left(\begin{smallmatrix} t^{i}&h(t)\\0&t^{j}\end{smallmatrix}\right) \in B_{2}^{+}(R)$ 
belongs to the Reidemeister class $\left[\left(\begin{smallmatrix} t^{x}&0\\0&t^{y}\end{smallmatrix}\right)\right]_{\phee_B}$ if and only if the following hold.
\begin{enumerate}
    \item There exist $k,\ell \in \Z$ such that $i = 2k+x$ and $j = 2\ell+y$, and
    \item there exists $f(t) \in \F_q[t,t^{-1}]$ such that 
    \begin{equation}
\label{eq:hf}h(t)=t^{\ell+y}f(t)-at^{2k+\ell+x}f(t^{-1}).
\end{equation}
\end{enumerate}
We now show that such an $f(t)$ always exists for any arbitrary (but fixed) $x,y,k,\ell \in \Z$ and $h(t) \in \F_q[t,t^{-1}]$.

Write $h(t)=\sum_{m\in\Z}h_m t^m$, where only finitely many $h_m\in\F_q$ are non-zero. 

For each $m \in \Z$, consider the number $\lambda_m=2k+2\ell+x+y-m \in \Z$. We observe that there can be at most one $m \in \Z$ for which the equality $m = \lambda_m$ holds. Now start by setting $f_m$ in the following cases. 
\[f_m=\begin{cases} (1-a)^{-1}h_m, &\text{ if } m=\lambda_m, \\ 
0, &\text{ if } h_m=h_{\lambda_m}=0.
\end{cases}\]
In case it holds both that $m \neq \lambda_m$ \emph{and} that $h_m$ and $h_{\lambda_m}$ are not {simultaneously} zero, we set each pair $(f_m,f_{\lambda_m}) \in \F_q^2$ to be a solution to the system 
\begin{align*} X-aY&=h_m\\
			  -aX+Y&=h_{\lambda_m}
\end{align*}
as in Lemma~\ref{lemmaa}. We remark that, since only finitely many of the coefficients $h_m \in \F_q$ are non-zero, there are only finitely many equations as above to be considered. 

By construction we have that only finitely many coefficients $f_m$ are non-zero, that is, $f(t)=\sum_{m \in \Z}f_mt^m \in \F_q[t,t^{-1}]$. Moreover,  
\begin{align*} t^{\ell+y}f(t)-at^{2k+\ell+x}f(t^{-1}) & = \sum_{m\in\Z} \left(f_mt^{m+\ell+y}-af_mt^{2k+\ell+x-m}\right)\\
& = \sum_{m\in\Z} (f_{m-\ell-y}-af_{2k+\ell+x-m})t^m\\
& = \sum_{m\in\Z} h_mt^m = h(t),
\end{align*}
hence equality~\bref{eq:hf} is satisfied for this choice of $f(t)$. 

The fact that such an $f(t) \in \F_q[t,t^{-1}]$ can always be constructed means that an arbitrary element $b=\left(\begin{smallmatrix} t^{i}&h(t)\\0&t^{j}\end{smallmatrix}\right) \in B_{2}^{+}(R)$ is contained in the Reidemeister class $\left[\left(\begin{smallmatrix} t^{x}&0\\0&t^{y}\end{smallmatrix}\right)\right]_{\phee_B}$ if and only if 
$i \equiv x \mod 2$ and $j \equiv y \mod 2$, without any further assumptions on $h(t)$. Therefore, 
\[\mathcal{R}(\phee_{B})=\left\{\left[\left(\begin{smallmatrix} 1&0\\0&1 \end{smallmatrix}\right)\right]_{\phee},\left[\left(\begin{smallmatrix} t&0\\0&1 \end{smallmatrix}\right)\right]_{\phee},\left[\left(\begin{smallmatrix} 1&0\\0&t \end{smallmatrix}\right)\right]_{\phee},\left[\left(\begin{smallmatrix} t&0\\0&t \end{smallmatrix}\right)\right]_{\phee}\right\},\]
as claimed.
\end{proof}

\section{Proof of Theorem \ref{thm:Lieapplication}, and final remarks} \label{sec:backtoarithmetic}

The promised theorem on $S$-arithmetic groups is now an obvious consequence of Propositions~\ref{pps:FptFptt-1OK} and~\ref{nori} together with known results from the literature.

\begin{proofof}{Theorem \ref{thm:Lieapplication}}
We define the ground field $\K_{n,p}$ and the $S$-arithmetic groups $\Gamma_{n,p} \leq \PB_n(\K_{n,p})$ case-by-case. 

If the chosen characteristic $p$ equals zero, set $\K_{n,p} = \Q(i)$ for all $n$, whose ring of integers is $\Z[i]$. 
For the arithmetic groups, take 
\[\Gamma_{2,0} = \mbU_2(\Z[i]) \rtimes_J C_3 \quad \text{ and } \quad  \Gamma_{n,0} = \PB_n(\Z[i]) \text{ for } n \geq 3,\]
where the generator of the cyclic part of $\Gamma_{2,0}$ acts on $\mbU_2(\Z[i]) \cong \Z^2$ via the matrix $J = \left( \begin{smallmatrix} -1 & 1 \\ -1 & 0 \end{smallmatrix} \right)$. The groups $\Gamma_{n,0} \leq \PB_n(\Q(i))$ with $n \geq 3$ are obviously arithmetic subgroups of $\PB_n$, while $\Gamma_{2,0}$ is an arithmetic subgroup of $\PB_2$ since it is commensurable with $\PB_2(\Z[i])$ via their common subgroup of finite index $\mbU_2(\Z[i])$. By Proposition~\ref{pps:FptFptt-1OK}, $\Gamma_{n,0}$ has property~$\Ri$ whenever $n \geq 4$. To see that $\RSpec(\Gamma_{2,0}) \neq \{1,\infty\}$ we draw from work of Dekimpe--Pennickx. More precisely, by (the proof of) \cite[Theorem~4.1]{KarelPenni}, the crystallographic group $\Gamma_{2,0} \cong \Z^2 \rtimes_J C_3$ admits an automorphism $\phee_D$ with finite Reidemeister number given by conjugation by the matrix $D = \left( \begin{smallmatrix} 0 & 1 \\ -1 & 1 \end{smallmatrix} \right)$. Since $D$ and $J$ commute, the map $\phee_D$ leaves the quotient $C_3$ invariant and thus $\phee_D$ has at least three twisted conjugacy classes. Thus $\RSpec(\Gamma_{2,0}) \neq \{1,\infty\}$. (The fact that all $\Gamma_{n,0}$ are finitely presented is well-known, but can also be seen directly from the fact that $\Gamma_{n,0}$ is virtually polycyclic.)

If the chosen characteristic $p$ is at least $5$, we take $\K_{n,p} = \F_p(t)$ for all $n$. Setting $\Gamma_{n,p} = \PB_n(\F_p[t])$ gives, by Propositions~\ref{nori}\bref{nori.i} and~\ref{pps:FptFptt-1OK}, the required properties for the Reidemeister spectra. Alternatively, take 
\[ \Gamma_{2,p} = \PBplus_2(\F_p[t,t^{-1}]) \quad \text{ and } \quad \Gamma_{n,p} = \PB_n(\F_p[t,t^{-1}]) \text{ for } n \geq 3, \] 
noticing that $\Gamma_{2,p}$ is commensurable with $\PB_2(\F_p[t,t^{-1}])$, and apply Propositions~\ref{nori}\bref{nori.ii} and~\ref{pps:FptFptt-1OK}. That the groups in the former (resp. latter) family are \emph{non}-finitely generated (resp. finitely generated) is also well-known; see, e.g., \cite[Theorem~1.3]{YuriSoluble}. 

In the remaining cases one needs to slightly vary the fields. If $p \in \{2,3\}$, set
\[ \K_{2,p} = \F_{p^2}(t) \quad \text{ and } \quad \K_{n,p} = \F_p(t) \text{ for } n \geq 3, \]
and take 
\[ \Gamma_{2,p} = \PBplus_2(\F_{p^2}[t,t^{-1}]) \quad \text{ and } \quad \Gamma_{n,p} = \PB_n(\F_p[t,t^{-1}]) \text{ for } n \geq 3. \]
The proof concludes as in the previous paragraph. 
\end{proofof}

The proof of Theorem~\ref{thm:Lieapplication} relied on the flexibility of the definition of $S$-arithmetic groups, which allows one to pass over to suitable commensurable groups. We observe, for instance, that the arithmetic subgroup $\PB_2(\Z[i]) \leq \PB_2(\Q(i))$ \emph{does} have property~$\Ri$, as can be seen from the classification of low-dimensional crystallographic groups with~$\Ri$ \cite{KarelPenni}. Similarly it is unknown whether, for example, $\PB_2(\F_5[t,t^{-1}])$ has property~$\Ri$, whence the descent to the subgroup of finite index $\PBplus_2(\F_5[t,t^{-1}])$.

On a similar token, to get a class-$2$ soluble, non-nilpotent, finitely generated $S$-arithmetic group $\Gamma$ with $\RSpec(\Gamma) \neq \{1,\infty\}$ and in characteristic two or three, we also needed to pass over to extensions of $\F_2(t)$ and $\F_3(t)$ since $\PBplus_2(\F_2[t,t^{-1}]) \cong C_2 \wr \Z$ and $\PBplus_2(\F_3[t,t^{-1}]) \cong C_3 \wr \Z$ have~$\Ri$ \cite{DacibergWongWreath}.

Taking these observations into account, and our lack of understanding of the case of class-$3$ soluble $S$-arithmetic groups, we conclude with the following. 

\begin{qst}
Let an arbitrary characteristic $p \geq 0$ be given. Does there exist a family of class-$n$ soluble linear algebraic groups $\{\mbG_n\}_{n\geq 2}$, defined over \emph{the same} global field $\K$ of characteristic $\carac(\K) = p$, together with a family $\{\Gamma_{n,p}\}_{n\geq 2}$ of class-$n$ soluble $S$-arithmetic subgroups of $\mbG_n$ such that 
\[ \RSpec(\Gamma_{2,p}) \neq \{1,\infty\} \neq \RSpec(\Gamma_{3,p}) \, \text{ but } \, \RSpec(\Gamma_{n,p}) = \{\infty\} \text{ for } n > 3? \]
\end{qst}


\section*{Acknowledgments}
PMLA was supported by the long term structural funding \emph{Methusalem grant} of the Flemish Government. YSR was partially supported by the German Research Foundation (DFG) through the Priority Program 2026 `Geometry at infinity'. 
Part of this project was inspired by a talk given by Peter Wong at the ICM 2018 Satellite Conference ``Geometric Group Theory'' in Campinas, Brazil. We thank him and Timur Nasybullov for helpful discussions. We are also indebted to Karel Dekimpe for valuable remarks, particularly for pointing us to a strengthening of Theorem~\ref{thm:Additive}.


 \def\cprime{$'$} \def\cprime{$'$}
 \providecommand{\bysame}{\leavevmode\hbox to3em{\hrulefill}\thinspace}
 \providecommand{\MR}{\relax\ifhmode\unskip\space\fi MR }

\printbibliography

\end{document}
% $ biblatex auxiliary file $
% $ biblatex bbl format version 3.2 $
% Do not modify the above lines!
%
% This is an auxiliary file used by the 'biblatex' package.
% This file may safely be deleted. It will be recreated as
% required.
%
\begingroup
\makeatletter
\@ifundefined{ver@biblatex.sty}
  {\@latex@error
     {Missing 'biblatex' package}
     {The bibliography requires the 'biblatex' package.}
      \aftergroup\endinput}
  {}
\endgroup

\datalist[entry]{nyt/global//global/global}
  \entry{BhuniaBose1}{article}{}
    \name{author}{2}{}{%
      {{hash=BS}{%
         family={Bhunia},
         familyi={B\bibinitperiod},
         given={S.},
         giveni={S\bibinitperiod},
      }}%
      {{hash=BA}{%
         family={Bose},
         familyi={B\bibinitperiod},
         given={A.},
         giveni={A\bibinitperiod},
      }}%
    }
    \strng{namehash}{BSBA1}
    \strng{fullhash}{BSBA1}
    \field{labelnamesource}{author}
    \field{labeltitlesource}{title}
    \field{sortinit}{B}
    \field{sortinithash}{B}
    \verb{doi}
    \verb 10.1007/s00031-020-09626-9
    \endverb
    \field{title}{Twisted conjugacy in linear algebraic groups}
    \field{journaltitle}{Transform. Groups}
    \field{year}{2020}
  \endentry

  \entry{BhuniaBose2}{article}{}
    \name{author}{2}{}{%
      {{hash=BS}{%
         family={Bhunia},
         familyi={B\bibinitperiod},
         given={S.},
         giveni={S\bibinitperiod},
      }}%
      {{hash=BA}{%
         family={Bose},
         familyi={B\bibinitperiod},
         given={A.},
         giveni={A\bibinitperiod},
      }}%
    }
    \strng{namehash}{BSBA1}
    \strng{fullhash}{BSBA1}
    \field{labelnamesource}{author}
    \field{labeltitlesource}{title}
    \field{sortinit}{B}
    \field{sortinithash}{B}
    \verb{doi}
    \verb 10.1016/j.jalgebra.2022.03.031
    \endverb
    \field{pages}{235\bibrangedash 259}
    \field{title}{Twisted conjugacy in linear algebraic groups {II}}
    \field{volume}{603}
    \field{journaltitle}{J. Algebra}
    \field{year}{2022}
  \endentry

  \entry{Bux04}{article}{}
    \name{author}{1}{}{%
      {{hash=BKU}{%
         family={Bux},
         familyi={B\bibinitperiod},
         given={K.-U.},
         giveni={K\bibinithyphendelim U\bibinitperiod},
      }}%
    }
    \strng{namehash}{BKU1}
    \strng{fullhash}{BKU1}
    \field{labelnamesource}{author}
    \field{labeltitlesource}{title}
    \field{sortinit}{B}
    \field{sortinithash}{B}
    \verb{doi}
    \verb 10.2140/gt.2004.8.611
    \endverb
    \field{pages}{611\bibrangedash 644}
    \field{title}{Finiteness properties of soluble arithmetic groups over
  global function fields}
    \field{volume}{8}
    \field{journaltitle}{Geom. Topol.}
    \field{year}{2004}
  \endentry

  \entry{KarelDacibergNil}{article}{}
    \name{author}{2}{}{%
      {{hash=DK}{%
         family={Dekimpe},
         familyi={D\bibinitperiod},
         given={K.},
         giveni={K\bibinitperiod},
      }}%
      {{hash=GD}{%
         family={Gon\c{c}alves},
         familyi={G\bibinitperiod},
         given={D.},
         giveni={D\bibinitperiod},
      }}%
    }
    \strng{namehash}{DKGD1}
    \strng{fullhash}{DKGD1}
    \field{labelnamesource}{author}
    \field{labeltitlesource}{title}
    \field{sortinit}{D}
    \field{sortinithash}{D}
    \verb{doi}
    \verb 10.1112/blms/bdu029
    \endverb
    \field{number}{4}
    \field{pages}{737\bibrangedash 746}
    \field{title}{The {$R_\infty$} property for free groups, free nilpotent
  groups and free solvable groups}
    \field{volume}{46}
    \field{journaltitle}{Bull. Lond. Math. Soc.}
    \field{year}{2014}
  \endentry

  \entry{KarelDacibergAbelian}{article}{}
    \name{author}{2}{}{%
      {{hash=DK}{%
         family={Dekimpe},
         familyi={D\bibinitperiod},
         given={K.},
         giveni={K\bibinitperiod},
      }}%
      {{hash=GD}{%
         family={Gon\c{c}alves},
         familyi={G\bibinitperiod},
         given={D.},
         giveni={D\bibinitperiod},
      }}%
    }
    \strng{namehash}{DKGD1}
    \strng{fullhash}{DKGD1}
    \field{labelnamesource}{author}
    \field{labeltitlesource}{title}
    \field{sortinit}{D}
    \field{sortinithash}{D}
    \verb{doi}
    \verb 10.12775/TMNA.2015.066
    \endverb
    \field{number}{2}
    \field{pages}{773\bibrangedash 784}
    \field{title}{The {$R_\infty$} property for abelian groups}
    \field{volume}{46}
    \field{journaltitle}{Topol. Methods Nonlinear Anal.}
    \field{year}{2015}
  \endentry

  \entry{KarelKaiserSam}{article}{}
    \name{author}{3}{}{%
      {{hash=DK}{%
         family={Dekimpe},
         familyi={D\bibinitperiod},
         given={K.},
         giveni={K\bibinitperiod},
      }}%
      {{hash=KT}{%
         family={Kaiser},
         familyi={K\bibinitperiod},
         given={T.},
         giveni={T\bibinitperiod},
      }}%
      {{hash=TS}{%
         family={Tertooy},
         familyi={T\bibinitperiod},
         given={S.},
         giveni={S\bibinitperiod},
      }}%
    }
    \strng{namehash}{DKKTTS1}
    \strng{fullhash}{DKKTTS1}
    \field{labelnamesource}{author}
    \field{labeltitlesource}{title}
    \field{sortinit}{D}
    \field{sortinithash}{D}
    \verb{doi}
    \verb 10.1016/j.jalgebra.2019.04.038
    \endverb
    \field{pages}{353\bibrangedash 375}
    \field{title}{The {R}eidemeister spectra of low dimensional
  crystallographic groups}
    \field{volume}{533}
    \field{journaltitle}{J. Algebra}
    \field{year}{2019}
  \endentry

  \entry{KarelPenni}{article}{}
    \name{author}{2}{}{%
      {{hash=DK}{%
         family={Dekimpe},
         familyi={D\bibinitperiod},
         given={K.},
         giveni={K\bibinitperiod},
      }}%
      {{hash=PP}{%
         family={Penninckx},
         familyi={P\bibinitperiod},
         given={P.},
         giveni={P\bibinitperiod},
      }}%
    }
    \strng{namehash}{DKPP1}
    \strng{fullhash}{DKPP1}
    \field{labelnamesource}{author}
    \field{labeltitlesource}{title}
    \field{sortinit}{D}
    \field{sortinithash}{D}
    \verb{doi}
    \verb 10.1007/s11784-011-0043-2
    \endverb
    \field{number}{2}
    \field{pages}{257\bibrangedash 283}
    \field{title}{The finiteness of the {R}eidemeister number of morphisms
  between almost-crystallographic groups}
    \field{volume}{9}
    \field{journaltitle}{J. Fixed Point Theory Appl.}
    \field{year}{2011}
  \endentry

  \entry{KarelSam}{article}{}
    \name{author}{2}{}{%
      {{hash=DK}{%
         family={Dekimpe},
         familyi={D\bibinitperiod},
         given={K.},
         giveni={K\bibinitperiod},
      }}%
      {{hash=TS}{%
         family={Tertooy},
         familyi={T\bibinitperiod},
         given={S.},
         giveni={S\bibinitperiod},
      }}%
    }
    \strng{namehash}{DKTS1}
    \strng{fullhash}{DKTS1}
    \field{labelnamesource}{author}
    \field{labeltitlesource}{title}
    \field{sortinit}{D}
    \field{sortinithash}{D}
    \verb{doi}
    \verb 10.1016/j.topol.2020.107565
    \endverb
    \field{pages}{Paper No. 107565, pp. 12}
    \field{title}{Algorithms for twisted conjugacy classes of
  polycyclic-by-finite groups}
    \field{volume}{293}
    \field{journaltitle}{Topology Appl.}
    \field{year}{2021}
  \endentry

  \entry{DekTerVandeBus}{article}{}
    \name{author}{3}{}{%
      {{hash=DK}{%
         family={Dekimpe},
         familyi={D\bibinitperiod},
         given={K.},
         giveni={K\bibinitperiod},
      }}%
      {{hash=TS}{%
         family={Tertooy},
         familyi={T\bibinitperiod},
         given={S.},
         giveni={S\bibinitperiod},
      }}%
      {{hash=VI}{%
         family={Van{ den B}ussche},
         familyi={V\bibinitperiod},
         given={I.},
         giveni={I\bibinitperiod},
      }}%
    }
    \strng{namehash}{DKTSVI1}
    \strng{fullhash}{DKTSVI1}
    \field{labelnamesource}{author}
    \field{labeltitlesource}{title}
    \field{sortinit}{D}
    \field{sortinithash}{D}
    \verb{doi}
    \verb 10.1080/00927872.2018.1435792
    \endverb
    \field{number}{9}
    \field{pages}{4090\bibrangedash 4103}
    \field{title}{Reidemeister zeta functions of low-dimensional
  almost-crystallographic groups are rational}
    \field{volume}{46}
    \field{journaltitle}{Comm. Algebra}
    \field{year}{2018}
  \endentry

  \entry{KarelSamIrisSolv}{article}{}
    \name{author}{3}{}{%
      {{hash=DK}{%
         family={Dekimpe},
         familyi={D\bibinitperiod},
         given={K.},
         giveni={K\bibinitperiod},
      }}%
      {{hash=TS}{%
         family={Tertooy},
         familyi={T\bibinitperiod},
         given={S.},
         giveni={S\bibinitperiod},
      }}%
      {{hash=VI}{%
         family={Van{ den B}ussche},
         familyi={V\bibinitperiod},
         given={I.},
         giveni={I\bibinitperiod},
      }}%
    }
    \strng{namehash}{DKTSVI1}
    \strng{fullhash}{DKTSVI1}
    \field{labelnamesource}{author}
    \field{labeltitlesource}{title}
    \field{sortinit}{D}
    \field{sortinithash}{D}
    \verb{doi}
    \verb 10.12775/tmna.2019.012
    \endverb
    \field{number}{2}
    \field{pages}{575\bibrangedash 601}
    \field{title}{Reidemeister spectra for solvamanifolds in low dimensions}
    \field{volume}{53}
    \field{journaltitle}{Topol. Methods Nonlinear Anal.}
    \field{year}{2019}
  \endentry

  \entry{FelshtynDacibergMetabelian}{article}{}
    \name{author}{2}{}{%
      {{hash=FA}{%
         family={Fel{'}shtyn},
         familyi={F\bibinitperiod},
         given={A.},
         giveni={A\bibinitperiod},
      }}%
      {{hash=GD}{%
         family={Gon\c{c}alves},
         familyi={G\bibinitperiod},
         given={D.},
         giveni={D\bibinitperiod},
      }}%
    }
    \strng{namehash}{FAGD1}
    \strng{fullhash}{FAGD1}
    \field{labelnamesource}{author}
    \field{labeltitlesource}{title}
    \field{sortinit}{F}
    \field{sortinithash}{F}
    \verb{doi}
    \verb 10.1142/S0218196711006297
    \endverb
    \field{number}{3}
    \field{pages}{505\bibrangedash 520}
    \field{title}{Reidemeister spectrum for metabelian groups of the form
  {$Q^n\rtimes\mathbb{Z}$} and {$\mathbb{Z}[1/p]^n\rtimes\mathbb{Z}$}, {$p$}
  prime}
    \field{volume}{21}
    \field{journaltitle}{Internat. J. Algebra Comput.}
    \field{year}{2011}
  \endentry

  \entry{FelshtynHill}{article}{}
    \name{author}{2}{}{%
      {{hash=FA}{%
         family={Fel'shtyn},
         familyi={F\bibinitperiod},
         given={A.},
         giveni={A\bibinitperiod},
      }}%
      {{hash=HR}{%
         family={Hill},
         familyi={H\bibinitperiod},
         given={R.},
         giveni={R\bibinitperiod},
      }}%
    }
    \strng{namehash}{FAHR1}
    \strng{fullhash}{FAHR1}
    \field{labelnamesource}{author}
    \field{labeltitlesource}{title}
    \field{sortinit}{F}
    \field{sortinithash}{F}
    \verb{doi}
    \verb 10.1007/BF00961408
    \endverb
    \field{number}{4}
    \field{pages}{367\bibrangedash 393}
    \field{title}{The {R}eidemeister zeta function with applications to
  {N}ielsen theory and a connection with {R}eidemeister torsion}
    \field{volume}{8}
    \field{journaltitle}{$K$-Theory}
    \field{year}{1994}
  \endentry

  \entry{FelshtynTroitskyCrelle}{article}{}
    \name{author}{2}{}{%
      {{hash=FA}{%
         family={Fel{'}shtyn},
         familyi={F\bibinitperiod},
         given={A.},
         giveni={A\bibinitperiod},
      }}%
      {{hash=TE}{%
         family={Troitsky},
         familyi={T\bibinitperiod},
         given={E.},
         giveni={E\bibinitperiod},
      }}%
    }
    \strng{namehash}{FATE1}
    \strng{fullhash}{FATE1}
    \field{labelnamesource}{author}
    \field{labeltitlesource}{title}
    \field{sortinit}{F}
    \field{sortinithash}{F}
    \verb{doi}
    \verb 10.1515/CRELLE.2007.097
    \endverb
    \field{pages}{193\bibrangedash 210}
    \field{title}{Twisted {B}urnside--{F}robenius theory for discrete groups}
    \field{volume}{613}
    \field{journaltitle}{J. reine angew. {M}ath.}
    \field{year}{2007}
  \endentry

  \entry{FelshtynZietek}{incollection}{}
    \name{author}{2}{}{%
      {{hash=FA}{%
         family={Fel{'}shtyn},
         familyi={F\bibinitperiod},
         given={A.},
         giveni={A\bibinitperiod},
      }}%
      {{hash=ZM}{%
         family={Zietek},
         familyi={Z\bibinitperiod},
         given={M.},
         giveni={M\bibinitperiod},
      }}%
    }
    \name{editor}{4}{}{%
      {{hash=MP}{%
         family={Moree},
         familyi={M\bibinitperiod},
         given={P.},
         giveni={P\bibinitperiod},
      }}%
      {{hash=PA}{%
         family={Pohl},
         familyi={P\bibinitperiod},
         given={A.},
         giveni={A\bibinitperiod},
      }}%
      {{hash=SL}{%
         family={Snoha},
         familyi={S\bibinitperiod},
         given={L.},
         giveni={L\bibinitperiod},
      }}%
      {{hash=WT}{%
         family={Ward},
         familyi={W\bibinitperiod},
         given={T.},
         giveni={T\bibinitperiod},
      }}%
    }
    \list{publisher}{1}{%
      {Amer. Math. Soc., Providence, RI}%
    }
    \strng{namehash}{FAZM1}
    \strng{fullhash}{FAZM1}
    \field{labelnamesource}{author}
    \field{labeltitlesource}{title}
    \field{sortinit}{F}
    \field{sortinithash}{F}
    \field{booktitle}{Dynamics: topology and numbers}
    \verb{doi}
    \verb 10.1090/conm/744/14979
    \endverb
    \field{pages}{57\bibrangedash 81}
    \field{series}{Contemp. Math.}
    \field{title}{Dynamical zeta functions of {R}eidemeister type and
  representations spaces}
    \field{volume}{744}
    \field{year}{2020}
  \endentry

  \entry{Daciberg}{article}{}
    \name{author}{1}{}{%
      {{hash=GD}{%
         family={Gon\c{c}alves},
         familyi={G\bibinitperiod},
         given={D.},
         giveni={D\bibinitperiod},
      }}%
    }
    \strng{namehash}{GD1}
    \strng{fullhash}{GD1}
    \field{labelnamesource}{author}
    \field{labeltitlesource}{title}
    \field{sortinit}{G}
    \field{sortinithash}{G}
    \verb{doi}
    \verb 10.1016/S0166-8641(97)00106-5
    \endverb
    \field{number}{3}
    \field{pages}{169\bibrangedash 186}
    \field{title}{Coincidence {R}eidemeister classes on nilmanifolds and
  nilpotent fibrations}
    \field{volume}{83}
    \field{journaltitle}{Topology Appl.}
    \field{year}{1998}
  \endentry

  \entry{DesiDaciberg}{article}{}
    \name{author}{2}{}{%
      {{hash=GD}{%
         family={Gon\c{c}alves},
         familyi={G\bibinitperiod},
         given={D.},
         giveni={D\bibinitperiod},
      }}%
      {{hash=KDH}{%
         family={Kochloukova},
         familyi={K\bibinitperiod},
         given={D.\bibnamedelima H.},
         giveni={D\bibinitperiod\bibinitdelim H\bibinitperiod},
      }}%
    }
    \strng{namehash}{GDKDH1}
    \strng{fullhash}{GDKDH1}
    \field{labelnamesource}{author}
    \field{labeltitlesource}{title}
    \field{sortinit}{G}
    \field{sortinithash}{G}
    \verb{doi}
    \verb 10.2140/pjm.2010.247.335
    \endverb
    \field{number}{2}
    \field{pages}{335\bibrangedash 352}
    \field{title}{Sigma theory and twisted conjugacy classes}
    \field{volume}{247}
    \field{journaltitle}{Pacific J. Math.}
    \field{year}{2010}
  \endentry

  \entry{DacibergWongWreath}{article}{}
    \name{author}{2}{}{%
      {{hash=GD}{%
         family={Gon\c{c}alves},
         familyi={G\bibinitperiod},
         given={D.},
         giveni={D\bibinitperiod},
      }}%
      {{hash=WP}{%
         family={Wong},
         familyi={W\bibinitperiod},
         given={P.},
         giveni={P\bibinitperiod},
      }}%
    }
    \strng{namehash}{GDWP1}
    \strng{fullhash}{GDWP1}
    \field{labelnamesource}{author}
    \field{labeltitlesource}{title}
    \field{sortinit}{G}
    \field{sortinithash}{G}
    \verb{doi}
    \verb 10.1142/S0218196706003219
    \endverb
    \field{pages}{875\bibrangedash 886}
    \field{title}{Twisted conjugacy classes in wreath products}
    \field{volume}{16}
    \field{journaltitle}{Internat. J. Algebra Comput.}
    \field{year}{2006}
  \endentry

  \entry{DacibergWongCrelle}{article}{}
    \name{author}{2}{}{%
      {{hash=GD}{%
         family={Gon\c{c}alves},
         familyi={G\bibinitperiod},
         given={D.},
         giveni={D\bibinitperiod},
      }}%
      {{hash=WP}{%
         family={Wong},
         familyi={W\bibinitperiod},
         given={P.},
         giveni={P\bibinitperiod},
      }}%
    }
    \strng{namehash}{GDWP1}
    \strng{fullhash}{GDWP1}
    \field{labelnamesource}{author}
    \field{labeltitlesource}{title}
    \field{sortinit}{G}
    \field{sortinithash}{G}
    \verb{doi}
    \verb 10.1515/CRELLE.2009.058
    \endverb
    \field{pages}{11\bibrangedash 27}
    \field{title}{Twisted conjugacy classes in nilpotent groups}
    \field{volume}{633}
    \field{journaltitle}{J. reine angew. {M}ath.}
    \field{year}{2009}
  \endentry

  \entry{HahnO'Meara}{book}{}
    \name{author}{2}{}{%
      {{hash=HAJ}{%
         family={Hahn},
         familyi={H\bibinitperiod},
         given={A.\bibnamedelima J.},
         giveni={A\bibinitperiod\bibinitdelim J\bibinitperiod},
      }}%
      {{hash=OOT}{%
         family={O'Meara},
         familyi={O\bibinitperiod},
         given={O.\bibnamedelima T.},
         giveni={O\bibinitperiod\bibinitdelim T\bibinitperiod},
      }}%
    }
    \list{publisher}{1}{%
      {Springer-Verlag, Berlin}%
    }
    \strng{namehash}{HAJOOT1}
    \strng{fullhash}{HAJOOT1}
    \field{labelnamesource}{author}
    \field{labeltitlesource}{title}
    \field{sortinit}{H}
    \field{sortinithash}{H}
    \verb{doi}
    \verb 10.1007/978-3-662-13152-7
    \endverb
    \field{pages}{xvi+576}
    \field{series}{Grundlehren der Mathematischen Wissenschaften}
    \field{title}{The classical groups and {$K$}-theory}
    \field{volume}{291}
    \field{year}{1989}
  \endentry

  \entry{HalasiPalfy}{article}{}
    \name{author}{2}{}{%
      {{hash=HZ}{%
         family={Halasi},
         familyi={H\bibinitperiod},
         given={Z.},
         giveni={Z\bibinitperiod},
      }}%
      {{hash=PPP}{%
         family={P{\'a}lfy},
         familyi={P\bibinitperiod},
         given={P.\bibnamedelima P.},
         giveni={P\bibinitperiod\bibinitdelim P\bibinitperiod},
      }}%
    }
    \strng{namehash}{HZPPP1}
    \strng{fullhash}{HZPPP1}
    \field{labelnamesource}{author}
    \field{labeltitlesource}{title}
    \field{sortinit}{H}
    \field{sortinithash}{H}
    \verb{doi}
    \verb 10.1515/jgt.2010.081
    \endverb
    \field{number}{6}
    \field{pages}{841\bibrangedash 854}
    \field{title}{The number of conjugacy classes in pattern groups is not a
  polynomial function}
    \field{volume}{14}
    \field{journaltitle}{J. Group Theory}
    \field{year}{2011}
  \endentry

  \entry{Higman}{article}{}
    \name{author}{1}{}{%
      {{hash=HG}{%
         family={Higman},
         familyi={H\bibinitperiod},
         given={G.},
         giveni={G\bibinitperiod},
      }}%
    }
    \strng{namehash}{HG1}
    \strng{fullhash}{HG1}
    \field{labelnamesource}{author}
    \field{labeltitlesource}{title}
    \field{sortinit}{H}
    \field{sortinithash}{H}
    \verb{doi}
    \verb 10.1112/plms/s3-10.1.24
    \endverb
    \field{pages}{24\bibrangedash 30}
    \field{title}{Enumerating {$p$}-groups. {I}. {I}nequalities}
    \field{volume}{10}
    \field{journaltitle}{Proc. London Math. Soc. (3)}
    \field{year}{1960}
  \endentry

  \entry{Jabara}{article}{}
    \name{author}{1}{}{%
      {{hash=JE}{%
         family={Jabara},
         familyi={J\bibinitperiod},
         given={E.},
         giveni={E\bibinitperiod},
      }}%
    }
    \strng{namehash}{JE1}
    \strng{fullhash}{JE1}
    \field{labelnamesource}{author}
    \field{labeltitlesource}{title}
    \field{sortinit}{J}
    \field{sortinithash}{J}
    \verb{doi}
    \verb 10.1016/j.jalgebra.2008.09.001
    \endverb
    \field{number}{10}
    \field{pages}{3671\bibrangedash 3679}
    \field{title}{Automorphisms with finite {R}eidemeister number in residually
  finite groups}
    \field{volume}{320}
    \field{journaltitle}{J. Algebra}
    \field{year}{2008}
  \endentry

  \entry{LevchukOriginal}{article}{}
    \name{author}{1}{}{%
      {{hash=LVM}{%
         family={Levchuk},
         familyi={L\bibinitperiod},
         given={V.\bibnamedelima M.},
         giveni={V\bibinitperiod\bibinitdelim M\bibinitperiod},
      }}%
    }
    \list{language}{1}{%
      {In Russian}%
    }
    \strng{namehash}{LVM1}
    \strng{fullhash}{LVM1}
    \field{labelnamesource}{author}
    \field{labeltitlesource}{title}
    \field{sortinit}{L}
    \field{sortinithash}{L}
    \field{number}{4}
    \field{pages}{64\bibrangedash 80}
    \field{title}{{Connections of unitriangular groups with certain rings. Ch.
  2. Automorphism groups}}
    \verb{url}
    \verb http://eudml.org/doc/62311
    \endverb
    \field{volume}{24}
    \field{journaltitle}{{Sib. Mat. Zh.}}
    \field{year}{1983}
  \endentry

  \entry{PaulaII}{article}{}
    \name{author}{1}{}{%
      {{hash=LPM}{%
         family={Lins{ }de{ A}raujo},
         familyi={L\bibinitperiod},
         given={P.\bibnamedelima M.},
         giveni={P\bibinitperiod\bibinitdelim M\bibinitperiod},
      }}%
    }
    \strng{namehash}{LPM1}
    \strng{fullhash}{LPM1}
    \field{labelnamesource}{author}
    \field{labeltitlesource}{title}
    \field{sortinit}{L}
    \field{sortinithash}{L}
    \verb{doi}
    \verb 10.1142/S0218196720500265
    \endverb
    \field{number}{5}
    \field{pages}{931\bibrangedash 975}
    \field{title}{Bivariate representation and conjugacy class zeta functions
  associated to unipotent group schemes, {II}: {G}roups of type {$F$}, {$G$},
  and {$H$}}
    \field{volume}{30}
    \field{journaltitle}{Internat. J. Algebra Comput.}
    \field{year}{2020}
  \endentry

  \entry{Brasuca0}{article}{}
    \name{author}{3}{}{%
      {{hash=LPM}{%
         family={Lins{ }de{ A}raujo},
         familyi={L\bibinitperiod},
         given={P.\bibnamedelima M.},
         giveni={P\bibinitperiod\bibinitdelim M\bibinitperiod},
      }}%
      {{hash=dOTAS}{%
         prefix={de},
         prefixi={d\bibinitperiod},
         family={Oliveira-Tosti},
         familyi={O\bibinithyphendelim T\bibinitperiod},
         given={A.\bibnamedelima S.},
         giveni={A\bibinitperiod\bibinitdelim S\bibinitperiod},
      }}%
      {{hash=SY}{%
         family={Santos{ R}ego},
         familyi={S\bibinitperiod},
         given={Y.},
         giveni={Y\bibinitperiod},
      }}%
    }
    \strng{namehash}{LPMOTASdSY1}
    \strng{fullhash}{LPMOTASdSY1}
    \field{labelnamesource}{author}
    \field{labeltitlesource}{title}
    \field{sortinit}{L}
    \field{sortinithash}{L}
    \verb{doi}
    \verb 10.1007/s10711-023-00790-2
    \endverb
    \field{pages}{no. 54}
    \field{title}{Thompson-like groups, {R}eidemeister numbers, and fixed
  points}
    \field{volume}{217}
    \field{journaltitle}{Geom. Dedicata}
    \field{year}{2023}
  \endentry

  \entry{Bn1}{unpublished}{}
    \name{author}{2}{}{%
      {{hash=LPM}{%
         family={Lins{ }de{ A}raujo},
         familyi={L\bibinitperiod},
         given={P.\bibnamedelima M.},
         giveni={P\bibinitperiod\bibinitdelim M\bibinitperiod},
      }}%
      {{hash=SY}{%
         family={Santos{ R}ego},
         familyi={S\bibinitperiod},
         given={Y.},
         giveni={Y\bibinitperiod},
      }}%
    }
    \strng{namehash}{LPMSY1}
    \strng{fullhash}{LPMSY1}
    \field{labelnamesource}{author}
    \field{labeltitlesource}{title}
    \field{sortinit}{L}
    \field{sortinithash}{L}
    \field{note}{Preprint, arXiv:2007.02988v2}
    \field{pages}{X}
    \field{title}{Twisted conjugacy in soluble arithmetic groups}
    \verb{url}
    \verb https://arxiv.org/abs/2007.02988v2
    \endverb
    \field{year}{2023}
  \endentry

  \entry{LitterickThomasCRgood}{article}{}
    \name{author}{2}{}{%
      {{hash=LAJ}{%
         family={Litterick},
         familyi={L\bibinitperiod},
         given={A.\bibnamedelima J.},
         giveni={A\bibinitperiod\bibinitdelim J\bibinitperiod},
      }}%
      {{hash=TAR}{%
         family={Thomas},
         familyi={T\bibinitperiod},
         given={A.\bibnamedelima R.},
         giveni={A\bibinitperiod\bibinitdelim R\bibinitperiod},
      }}%
    }
    \strng{namehash}{LAJTAR1}
    \strng{fullhash}{LAJTAR1}
    \field{labelnamesource}{author}
    \field{labeltitlesource}{title}
    \field{sortinit}{L}
    \field{sortinithash}{L}
    \verb{doi}
    \verb 10.1090/tran/7085
    \endverb
    \field{number}{8}
    \field{pages}{5279\bibrangedash 5340}
    \field{title}{Complete reducibility in good characteristic}
    \field{volume}{370}
    \field{journaltitle}{Trans. Amer. Math. Soc.}
    \field{year}{2018}
  \endentry

  \entry{Margulis}{book}{}
    \name{author}{1}{}{%
      {{hash=MGA}{%
         family={Margulis},
         familyi={M\bibinitperiod},
         given={G.\bibnamedelima A.},
         giveni={G\bibinitperiod\bibinitdelim A\bibinitperiod},
      }}%
    }
    \list{publisher}{1}{%
      {Springer-Verlag, Berlin}%
    }
    \strng{namehash}{MGA1}
    \strng{fullhash}{MGA1}
    \field{labelnamesource}{author}
    \field{labeltitlesource}{title}
    \field{sortinit}{M}
    \field{sortinithash}{M}
    \field{isbn}{978-3-642-05721-2}
    \field{pages}{x+388}
    \field{series}{Ergebnisse der Mathematik und ihrer Grenzgebiete (3)}
    \field{title}{Discrete subgroups of semisimple {L}ie groups}
    \field{volume}{17}
    \field{year}{1991}
  \endentry

  \entry{TimurUniTri}{article}{}
    \name{author}{1}{}{%
      {{hash=NT}{%
         family={Nasybullov},
         familyi={N\bibinitperiod},
         given={T.},
         giveni={T\bibinitperiod},
      }}%
    }
    \strng{namehash}{NT1}
    \strng{fullhash}{NT1}
    \field{labelnamesource}{author}
    \field{labeltitlesource}{title}
    \field{sortinit}{N}
    \field{sortinithash}{N}
    \verb{doi}
    \verb 10.1515/jgth-2018-0127
    \endverb
    \field{number}{2}
    \field{pages}{253\bibrangedash 266}
    \field{title}{Twisted conjugacy classes in unitriangular groups}
    \field{volume}{22}
    \field{journaltitle}{J. Group Theory}
    \field{year}{2019}
  \endentry

  \entry{PlatonovRapinchuk}{book}{}
    \name{author}{2}{}{%
      {{hash=PV}{%
         family={Platonov},
         familyi={P\bibinitperiod},
         given={V.},
         giveni={V\bibinitperiod},
      }}%
      {{hash=RA}{%
         family={Rapinchuk},
         familyi={R\bibinitperiod},
         given={A.},
         giveni={A\bibinitperiod},
      }}%
    }
    \list{publisher}{1}{%
      {Academic Press Inc., Boston, MA}%
    }
    \strng{namehash}{PVRA1}
    \strng{fullhash}{PVRA1}
    \field{labelnamesource}{author}
    \field{labeltitlesource}{title}
    \field{sortinit}{P}
    \field{sortinithash}{P}
    \field{isbn}{0125581807}
    \field{pages}{xii+614}
    \field{series}{Pure and Applied Mathematics}
    \field{title}{Algebraic groups and number theory}
    \field{volume}{139}
    \field{year}{1994}
  \endentry

  \entry{Romankov}{article}{}
    \name{author}{1}{}{%
      {{hash=RV}{%
         family={Roman{'}kov},
         familyi={R\bibinitperiod},
         given={V.},
         giveni={V\bibinitperiod},
      }}%
    }
    \strng{namehash}{RV1}
    \strng{fullhash}{RV1}
    \field{labelnamesource}{author}
    \field{labeltitlesource}{title}
    \field{sortinit}{R}
    \field{sortinithash}{R}
    \verb{doi}
    \verb 10.1016/j.jpaa.2010.06.015
    \endverb
    \field{number}{4}
    \field{pages}{664\bibrangedash 671}
    \field{title}{Twisted conjugacy classes in nilpotent groups}
    \field{volume}{215}
    \field{journaltitle}{J. Pure Appl. Algebra}
    \field{year}{2011}
  \endentry

  \entry{ask}{article}{}
    \name{author}{1}{}{%
      {{hash=RT}{%
         family={Rossmann},
         familyi={R\bibinitperiod},
         given={T.},
         giveni={T\bibinitperiod},
      }}%
    }
    \strng{namehash}{RT1}
    \strng{fullhash}{RT1}
    \field{labelnamesource}{author}
    \field{labeltitlesource}{title}
    \field{sortinit}{R}
    \field{sortinithash}{R}
    \verb{doi}
    \verb 10.1112/plms.12159
    \endverb
    \field{number}{3}
    \field{pages}{574\bibrangedash 616}
    \field{title}{The average size of the kernel of a matrix and orbits of
  linear groups}
    \field{volume}{117}
    \field{journaltitle}{Proc. Lond. Math. Soc. (3)}
    \field{year}{2018}
  \endentry

  \entry{RoVo19}{article}{}
    \name{author}{2}{}{%
      {{hash=RT}{%
         family={Rossmann},
         familyi={R\bibinitperiod},
         given={T.},
         giveni={T\bibinitperiod},
      }}%
      {{hash=VC}{%
         family={Voll},
         familyi={V\bibinitperiod},
         given={C.},
         giveni={C\bibinitperiod},
      }}%
    }
    \strng{namehash}{RTVC1}
    \strng{fullhash}{RTVC1}
    \field{labelnamesource}{author}
    \field{labeltitlesource}{title}
    \field{sortinit}{R}
    \field{sortinithash}{R}
    \field{note}{To appear}
    \field{title}{Groups, graphs, and hypergraphs: average sizes of kernels of
  generic matrices with support constraints}
    \verb{url}
    \verb https://arxiv.org/abs/1908.09589
    \endverb
    \field{journaltitle}{Mem. Amer. Math. Soc.}
    \field{year}{2022}
  \endentry

  \entry{YuriSoluble}{article}{}
    \name{author}{1}{}{%
      {{hash=SY}{%
         family={Santos{ R}ego},
         familyi={S\bibinitperiod},
         given={Y.},
         giveni={Y\bibinitperiod},
      }}%
    }
    \strng{namehash}{SY1}
    \strng{fullhash}{SY1}
    \field{labelnamesource}{author}
    \field{labeltitlesource}{title}
    \field{sortinit}{S}
    \field{sortinithash}{S}
    \verb{doi}
    \verb 10.4153/S0008414X21000213
    \endverb
    \field{number}{5}
    \field{pages}{1209\bibrangedash 1243}
    \field{title}{On the finiteness length of some soluble linear groups}
    \field{volume}{74}
    \field{journaltitle}{Canad. J. Math.}
    \field{year}{2022}
  \endentry

  \entry{duSau05}{article}{}
    \name{author}{1}{}{%
      {{hash=dSMPF}{%
         prefix={du},
         prefixi={d\bibinitperiod},
         family={Sautoy},
         familyi={S\bibinitperiod},
         given={M.\bibnamedelima P.\bibnamedelima F.},
         giveni={M\bibinitperiod\bibinitdelim P\bibinitperiod\bibinitdelim
  F\bibinitperiod},
      }}%
    }
    \strng{namehash}{SMPFd1}
    \strng{fullhash}{SMPFd1}
    \field{labelnamesource}{author}
    \field{labeltitlesource}{title}
    \field{sortinit}{S}
    \field{sortinithash}{S}
    \verb{doi}
    \verb 10.1112/S0024609304003637
    \endverb
    \field{number}{1}
    \field{pages}{37\bibrangedash 44}
    \field{title}{Counting conjugacy classes}
    \field{volume}{37}
    \field{journaltitle}{Bull. London Math. Soc.}
    \field{year}{2005}
  \endentry

  \entry{Silvester}{article}{}
    \name{author}{1}{}{%
      {{hash=SJR}{%
         family={Silvester},
         familyi={S\bibinitperiod},
         given={J.\bibnamedelima R.},
         giveni={J\bibinitperiod\bibinitdelim R\bibinitperiod},
      }}%
    }
    \strng{namehash}{SJR1}
    \strng{fullhash}{SJR1}
    \field{labelnamesource}{author}
    \field{labeltitlesource}{title}
    \field{sortinit}{S}
    \field{sortinithash}{S}
    \verb{doi}
    \verb 10.1112/plms/s3-26.1.35
    \endverb
    \field{pages}{35\bibrangedash 56}
    \field{title}{On the {$K_{2}$} of a free associative algebra}
    \field{volume}{26}
    \field{journaltitle}{Proc. London Math. Soc. (3)}
    \field{year}{1973}
  \endentry

  \entry{SteinbergEndo}{article}{}
    \name{author}{1}{}{%
      {{hash=SR}{%
         family={Steinberg},
         familyi={S\bibinitperiod},
         given={R.},
         giveni={R\bibinitperiod},
      }}%
    }
    \strng{namehash}{SR1}
    \strng{fullhash}{SR1}
    \field{labelnamesource}{author}
    \field{labeltitlesource}{title}
    \field{sortinit}{S}
    \field{sortinithash}{S}
    \verb{doi}
    \verb 10.1090/memo/0080
    \endverb
    \field{number}{80}
    \field{pages}{108}
    \field{title}{Endomorphisms of linear algebraic groups}
    \field{journaltitle}{Mem. Amer. Math. Soc.}
    \field{year}{1968}
  \endentry

  \entry{Steinberg}{book}{}
    \name{author}{1}{}{%
      {{hash=SR}{%
         family={Steinberg},
         familyi={S\bibinitperiod},
         given={R.},
         giveni={R\bibinitperiod},
      }}%
    }
    \list{publisher}{1}{%
      {American Mathematical Society, Providence, RI}%
    }
    \strng{namehash}{SR1}
    \strng{fullhash}{SR1}
    \field{labelnamesource}{author}
    \field{labeltitlesource}{title}
    \field{sortinit}{S}
    \field{sortinithash}{S}
    \verb{doi}
    \verb 10.1090/ulect/066
    \endverb
    \field{pages}{xi+160}
    \field{series}{University Lecture Series}
    \field{title}{Lectures on {C}hevalley groups}
    \field{volume}{66}
    \field{year}{2016}
  \endentry

  \entry{TabackWong0}{article}{}
    \name{author}{2}{}{%
      {{hash=TJ}{%
         family={Taback},
         familyi={T\bibinitperiod},
         given={J.},
         giveni={J\bibinitperiod},
      }}%
      {{hash=WP}{%
         family={Wong},
         familyi={W\bibinitperiod},
         given={P.},
         giveni={P\bibinitperiod},
      }}%
    }
    \strng{namehash}{TJWP1}
    \strng{fullhash}{TJWP1}
    \field{labelnamesource}{author}
    \field{labeltitlesource}{title}
    \field{sortinit}{T}
    \field{sortinithash}{T}
    \verb{doi}
    \verb 10.1112/jlms/jdm024
    \endverb
    \field{number}{3}
    \field{pages}{705\bibrangedash 717}
    \field{title}{Twisted conjugacy and quasi-isometry invariance for
  generalized solvable {B}aumslag--{S}olitar groups}
    \field{volume}{75}
    \field{journaltitle}{J. Lond. Math. Soc. (2)}
    \field{year}{2007}
  \endentry

  \entry{Tiemeyer}{article}{}
    \name{author}{1}{}{%
      {{hash=TA}{%
         family={Tiemeyer},
         familyi={T\bibinitperiod},
         given={A.},
         giveni={A\bibinitperiod},
      }}%
    }
    \strng{namehash}{TA1}
    \strng{fullhash}{TA1}
    \field{labelnamesource}{author}
    \field{labeltitlesource}{title}
    \field{sortinit}{T}
    \field{sortinithash}{T}
    \verb{doi}
    \verb 10.1007/BF01235942
    \endverb
    \field{number}{2}
    \field{pages}{215\bibrangedash 223}
    \field{title}{A local-global principle for finiteness properties of
  {$S$}-arithmetic groups over number fields}
    \field{volume}{2}
    \field{journaltitle}{Transform. Groups}
    \field{year}{1997}
  \endentry

  \entry{VLA}{article}{}
    \name{author}{2}{}{%
      {{hash=VLA}{%
         family={Vera-L\'opez},
         familyi={V\bibinithyphendelim L\bibinitperiod},
         given={A.},
         giveni={A\bibinitperiod},
      }}%
      {{hash=AJM}{%
         family={Arregi},
         familyi={A\bibinitperiod},
         given={J.\bibnamedelima M.},
         giveni={J\bibinitperiod\bibinitdelim M\bibinitperiod},
      }}%
    }
    \strng{namehash}{VLAAJM1}
    \strng{fullhash}{VLAAJM1}
    \field{labelnamesource}{author}
    \field{labeltitlesource}{title}
    \field{sortinit}{V}
    \field{sortinithash}{V}
    \verb{doi}
    \verb 10.1016/S0024-3795(03)00371-9
    \endverb
    \field{pages}{85\bibrangedash 124}
    \field{title}{Conjugacy classes in unitriangular matrices}
    \field{volume}{370}
    \field{journaltitle}{Linear Algebra Appl.}
    \field{year}{2003}
  \endentry

  \entry{WongCrelle}{article}{}
    \name{author}{1}{}{%
      {{hash=WP}{%
         family={Wong},
         familyi={W\bibinitperiod},
         given={P.},
         giveni={P\bibinitperiod},
      }}%
    }
    \strng{namehash}{WP1}
    \strng{fullhash}{WP1}
    \field{labelnamesource}{author}
    \field{labeltitlesource}{title}
    \field{sortinit}{W}
    \field{sortinithash}{W}
    \verb{doi}
    \verb 10.1515/crll.2000.056
    \endverb
    \field{pages}{185\bibrangedash 204}
    \field{title}{Reidemeister number, {H}irsch rank, coincidences on
  polycyclic groups and solvmanifolds}
    \field{volume}{524}
    \field{journaltitle}{J. reine angew. {M}ath.}
    \field{year}{2000}
  \endentry
\enddatalist
\endinput
